% Gerado por tools/tei2latex.py a partir de revisado.xml. Regenerar, não editar.

%%>>> parte-1

%%>>> cap-4

%%>>> art-20

%%>>> t143

\emph{Z} apprehende o ente, emquanto é dotado de uma certa \emph{essencia} \pg{228}  (pois o objecto \emph{proprio} da \emph{intelligencia} é a \emph{essencia}), --- ao passo que a \emph{vontade} tende para o objecto, emquanto é dotado de \emph{existencia real} (pois um objecto, emquanto \emph{existe} na realidade, é bem, e por isso é termo da \emph{vontade}). Ora o ente considerado \pg{229}  na sua \emph{essencia} é mais abstracto e universal que o ente considerado na sua \emph{existencia real}; pois que a \emph{essencia} não inclue a \emph{existencia real}, mas esta inclue aquella. Logo a \emph{intelligencia} é mais nobre que a \emph{vontade} \footnote{A intelligencia, \emph{absolutamente} fallando, é mais nobre que a vontade; porque o objecto \emph{formal} d'aquella é mais universal que o d'esta.\par Mas, se attendermos aos objectos \emph{materiaes} das duas faculdades, umas vezes a intelligencia é superior á vontade, e outras esta é superior áquella.\par Quando o objecto é inferior ao espirito humano, a intelligencia é superior á vontade; assim conhecer uma pedra é mais perfeito do que amal-a.\par Quando, pelo contrario, o objecto é superior ao espirito humano, então a vontade é superior á intelligencia; assim o amor, que nos une a Deus, é mais perfeito e mais nobre do que o simples conhecimento do Ente Infinito. A razão é clara. Na primeira hypothese, o objecto, que é inferior ao espirito humano, recebe da intelligencia uma forma mais nobre do que a que existe no proprio objecto e que attrahe a vontade; assim o ser ideal, que a pedra adquire no nosso espirito, é mais nobre que o ser material que é proprio d'ella: --- como, pelo contrario, o objecto, que é superior ao nosso espirito, existe em si mesmo de um modo mais nobre do que existe na nossa intelligencia; assim a Bondade de \textsc{Deus}, que é o objecto do nosso amor, é infinitamente mais nobre e elevada na Essencia Divina, do que no nosso espirito, que a conhece (Cf. \emph{Sum. Theol.,} p. I, q. 82, a. 3).\par Concluimos que, \emph{absolutamente} fallando, a intelligencia é mais nobre que a vontade; porque as faculdades especificam-se pelos seus objectos \emph{formaes}, e não pelos seus objectos \emph{materiaes}.\par Todavia, se a intelligencia, pela sua elevação e extensão, é mais nobre que a vontade, esta faculdade é mais intima, mais central, mais profunda e por isso mais secreta e mysteriosa do que aquella. D'ahi a expressão commum a quasi todas as linguas: eu te amo \emph{do fundo do coração}. Por isso é mais facil conhecer os pensamentos da intelligencia, do que os affectos do coração. \emph{Larochefoucauld} dizia: « Nem todos os que conhecem o seu espirito, conhecem o seu coração ».}.

\part*{SECÇÃO TERCEIRA\\[0.3em]\large OPERAÇÕES DO HOMEM}
\addcontentsline{toc}{part}{SECÇÃO TERCEIRA --- OPERAÇÕES DO HOMEM}

%%>>> t144

\pg{230} \tese{144}{Divisão da secção terceira.}
Das faculdades emanam as operações, como as faculdades emanam da essencia. Por isso, tendo considerado a \emph{essencia} e as \emph{faculdades} do homem, devemos agora estudar as suas \emph{operações}. --- Dividimos a secção \emph{terceira} em \emph{dois} capitulos. No \emph{primeiro} capitulo tratamos das operações \emph{sensitivas} do homem; no \emph{segundo}, das suas operações \emph{intellectuaes} \footnote{Não tratamos aqui de proposito das operações \emph{vegetativas} do homem, não só porque no Artigo II do Capitulo I da Secção primeira da \emph{Anthropologia} fallámos das \emph{funcções} da vida \emph{vegetativa}, mas tambem porque o assumpto não é tão importante, que exija maior desenvolvimento.\par As operações \emph{vegetativas} do homem, como tambem de todos os seres vivos mais perfeitos, são tres: \emph{nutrição}, \emph{crescimento} e \emph{reproducção}.\par Todavia diremos alguma coisa acerca do \emph{principio remoto} d'essas operações.\par Alguns, como \emph{Descartes}, dizem que as operações \emph{vegetativas} no homem derivam de um principio distincto e diverso da alma humana.\par Outros, como \emph{Stahl}, sustentam que as operações \emph{vegetativas} no homem derivam da alma \emph{racional} como \emph{tal}, e que por isso não são \emph{naturaes}, como nos outros seres destituidos de razão, mas são \emph{racionaes}.\par Outros, como \emph{Santo Thomaz}, seguindo uma doutrina media entre aquelles dois extremos, ensinam que as operações \emph{vegetativas} no homem derivam da alma \emph{racional}, não emquanto é \emph{racional}, mas emquanto \emph{virtualmente} contém o grau de vida \emph{vegetativa}.\par Façamos uma breve critica d'estas opiniões.\par a) A opinião de \emph{Descartes} é falsa. Com effeito, a alma \emph{racional}, como demonstrámos no \emph{Capitulo} III da Secção I d'este tratado, é a \emph{unica} forma existente em nós e o principio de todas as nossas operações.\par \emph{b}) A opinião de \emph{Stahl} tambem é falsa. Por quanto, se as operações \emph{vegetativas} derivassem de nossa alma, emquanto \emph{racional}, emanariam das faculdades \emph{intellectuaes}, e por isso seriam \emph{espirituaes}. Ora as operações \emph{vegetativas} não são \emph{espirituaes}, mas \emph{materiaes}, pois que o seu objecto é o corpo, que é material. Logo não derivam da alma, emquanto \emph{racional}. --- De mais, se as operações \emph{vegetativas} derivassem da alma, emquanto \emph{racional}, esta poderia dirigil-as e ter consciencia d'ellas, porque as operações \emph{racionaes} estão sujeitas á vontade e são percebidas pela \emph{consciencia}. Ora a alma não tem poder sobre as operações \emph{vegetativas}, nem tem consciencia d'ellas.\par Nem vale dizer com \emph{Stahl} que as paixões, e tambem as operações intellectuaes influem nas funcções \emph{vegetativas}. Esse influxo não significa que todas as operações humanas derivem da alma, emquanto \emph{racional}, mas indica\pgn{231}que em nós existe uma só alma de que emanam varias faculdades, e que essas faculdades exercem um mutuo influxo entre si.\par \emph{c}) A opinião de \emph{Santo Thomaz} é verdadeira. Com effeito, como provámos contra \emph{Descartes}, as operações \emph{vegetativas} derivam da nossa alma; --- e, como demonstrámos contra \emph{Stahl}, aquellas operações derivam da nossa alma, não emquanto é \emph{racional}, mas emquanto é \emph{virtualmente} vegetativa, isto é, emquanto tem a virtude ou força de produzir as operações \emph{vegetativas}.}.

\textbf{OPERAÇÕES SENSITIVAS DO HOMEM}

\textsc{SUMMARIO}: --- Percepção sensitiva, seu objecto e divisão. --- Percepção sensitiva

externa. --- Percepção sensitiva interna. --- Appetite sensitivo, sua

divisão e paixões.

%%>>> art-21

\section*{ARTIGO I.}
\addcontentsline{toc}{section}{ARTIGO I.}

\pg{231} \textbf{Percepção sensitiva}, \textbf{seu objecto} e \textbf{divisão}.

%%>>> t145

\tese{145}{Operações sensitivas do homem.}
As operações \emph{sensitivas} do homem são as que emanam \emph{immediatamente} das faculdades \emph{sensitivas}, de que elle é dotado. --- E, como as especies das faculdades \emph{sensitivas} são a \emph{perceptiva}, a \emph{appetitiva} e a \emph{locomotiva} com a \emph{vocal}, as operações \emph{sensitivas} do homem dividem-se em \emph{perceptivas}, \emph{appetitivas}, \emph{locomotivas} e \emph{vocaes}.

Neste e nos artigos seguintes trataremos das operações, que emanam das faculdades \emph{sensitivo-perceptivas} e \emph{sensitivo-appetitivas}, e que são a \emph{percepção sensitiva} e o \emph{appetite sensitivo} \footnote{Não reservamos artigo algum para o tratado das operações \emph{locomotivas} e \emph{vocaes}, porque são pouco importantes para o nosso scopo, e porque já falámos sobre o assumpto, quando expozemos a doutrina acerca das faculdades correspondentes, que são a \emph{locomotiva} e a \emph{vocal}.}.

%%>>> t146

\tese{146}{Percepção sensitiva.}
\emph{Percepção sensitiva} é uma operação pela qual o homem conhece \emph{directamente} os entes \emph{materiaes} na sua \emph{singularidade}. Assim a \emph{visão} é uma percepção \emph{sensitiva}, pois é por ella que o homem percebe directamente as \emph{côres}, que são objectos materiaes e concretos.

Como já dissemos, a operação, pela qual o homem conhece \emph{directamente} os objectos \emph{immateriaes}, ou os objectos \emph{materiaes} na sua \emph{universalidade}, \emph{é intellectual} \footnote{Dissemos que o homem, pela percepção \emph{sensitiva}, conhece \emph{directamente} os objectos \emph{materiaes} na sua \emph{singularidade}; pois que, embora elle, pela percepção \emph{intellectual}, conheça tambem os objectos \emph{materiaes} na sua \emph{singularidade}, este conhecimento é \emph{indirecto}, como explicaremos.\par O tratado do conhecimento \emph{sensitivo}, como tambem o das faculdades\pgn{232}\emph{sensitivas}, é muito util e necessario. Este conhecimento serve não só para o bem material do ser sensitivo, mas, no homem, tambem serve para o conhecimento \emph{intellectual}, emquanto a intelligencia abstrahe dos phantasmas, que a imaginação apresenta, as especies intelligiveis. Diz S. Thomaz: « Cognitio \emph{sensitiva} ordinatur ad duo: uno enim modo, tam in hominibus quam in aliis animalibus, ordinatur ad corporis sustentationem, quia per hujusmodi cognitionem homines et alia animalia vitant nociva et conquirunt ea, quae sunt necessaria ad corporis sustentationem: alio modo specialiter in homine ordinatur ad cognitionem intellectivam, vel speculativam vel practicam ». (\emph{Sum. Theol.}, 2-2, q. 167, a. 2).}.

147. \textbf{Objecto da percepção sensitiva.} --- O objecto da \emph{percepção sensitiva} é o ente \emph{material}, emquanto \emph{tal}. Por quanto, todo o objecto percebido deve ser proporcionado á faculdade perceptiva, de modo que, se a faculdade fôr material, o objecto tambem deve ser material. Ora a faculdade \emph{perceptivo-sensitiva é material}; pois exerce-se por meio de um orgão corporeo, que é material. Logo o objecto da faculdade \emph{perceptivo-sensitiva}, e por isso da \emph{percepção sensitiva}, é o ente \emph{material} \footnote{O objecto da percepção \emph{sensitiva} é constituido pelos seres corporeos, considerados na sua \emph{singularidade}, e por isso dotados de \emph{quantidade}. Diz S. Thomaz: « Sendo o sentido uma energia existente na quantidade, não pode ter por objecto senão um ente dotado de quantidade: porque deve haver proporção entre a faculdade e o seu objecto, entre o paciente e o agente: \emph{Sensus est virtus in magnitudine}, \emph{cum sit actus organi corporei}; \emph{et ideo non potest pati nisi ab habente magnitudinem}. \emph{Activum enim debet esse proportionatum passivo} » (\emph{De sensu et sensato}, l. 15). --- Todavia esses seres não são attingidos pela percepção \emph{sensitiva} em todos os seus aspectos ou attributos, mas só nas suas qualidades \emph{sensiveis} e \emph{singulares}. D'este modo os seres corporeos e singulares formam o objecto \emph{material} da percepção \emph{sensitiva}, e as qualidades sensiveis e singulares dos mesmos seres constituem o seu objecto \emph{formal}.\par Repetimos que as qualidades \emph{sensiveis}, objecto \emph{formal} da percepção \emph{sensitiva}, são \emph{individuas} ou \emph{concretas}, como se encontram na natureza corporea. Porquanto, as qualidades \emph{sensiveis} podem ser percebidas por uma faculdade cognitiva --- ou emquanto são abstrahidas do objecto material e corporeo, e são transportadas, não \emph{real} mas \emph{intencionalmente}, para a faculdade sensitiva por meio de uma imagem, que as represente na sua \emph{condição individua} e \emph{concreta}, --- ou emquanto são abstrahidas tambem da \emph{condição concreta} e \emph{singular}, de modo que só seja representada a sua natureza \emph{generica} ou \emph{especifica}. Consideradas naquelle \emph{primeiro} grau de abstracção, as qualidades \emph{sensiveis} são o objecto \emph{proprio} das faculdades \emph{sensitivas}; consideradas no \emph{segundo} grau, são o objecto \emph{proprio} da \emph{intelligencia}, que só considera as naturezas \emph{abstractas}, tanto dos entes, como das qualidades. Diz S. Thomaz: « \emph{Sensibilium qualitatum naturas} cognoscere non est sensus, sed intellectus ». (\emph{Sum. Theol.}, p. I, q. 78, a. 3).\par Os sentidos não só não percebem a natureza das qualidades \emph{sensiveis}, mas nem apprehendem umas \emph{relações}, que a intelligencia, pela abstracção que lhe é propria, reconhece nos objectos corporeos, e cuja percepção se attribue tambem, mas \emph{impropriamente}, á faculdade sensitiva. Taes são a\pgn{233}\emph{ordem}, a \emph{harmonia}, a \emph{proporção}, a \emph{belleza}. --- Dissemos --- \emph{impropriamente}; porque os sentidos, se reproduzem e conhecem um objecto sensivel \emph{ordenado}, \emph{harmonico}, \emph{proporcionado}, \emph{bello}, não reproduzem nem conhecem as imagens abstractas de \emph{belleza}, de \emph{ordem}, de \emph{harmonia}, de \emph{proporção}, e por isso não podem \emph{conhecer} a \emph{belleza} do que é bello, nem a \emph{ordem} do que é ordenado, nem a \emph{proporção} do que é proporcionado, nem a \emph{harmonia} do que é harmonico. Se o \emph{sentido} podesse receber a imagem \emph{abstracta} d'essas relações e conhecer a \emph{ordem}, etc., não seria faculdade \emph{sensitiva}, mas \emph{intellectual}.\par Devemos fazer uma observação importante. As \emph{qualidades}, que constituem o objecto dos nossos sentidos, não são todas as qualidades \emph{corporeas}, mas são apenas as que formam, como vimos na \emph{Ontologia}, a \emph{terceira} especie e se chamam \emph{qualidades passiveis}. Na verdade, embora \emph{todas} as qualidades produzam algum effeito, porque todas são principios \emph{formaes}, todavia só as da \emph{terceira} especie têm a virtude de produzir mudança ou alteração nos corpos, animados ou inanimados (\emph{Sum. Theol.}, p. I, q. 78, a. 3 ad 1). Taes são a \emph{luz}, a \emph{côr}, o \emph{calor}, o \emph{cheiro}, a \emph{humidade}, etc. e, em geral, todas as \emph{actividades} da substancia corporea, capazes de produzir uma alteração nos corpos. Diz S. Thomaz: « Unaquaeque praedictarum qualitatum est nata agere in sensum: in hoc enim propria ratio uniuscujusque consistit ut moveat sensum, sicut ad rationem coloris pertinet quod possit movere visum » (\emph{De sensu et sensato}, l. 15). --- E tendo, como dissemos, por sujeito a \emph{quantidade} da substancia corporea, estas qualidades não podem actuar senão nos \emph{sujeitos} extensos, e por \emph{meios} extensos, e por \emph{modos} subordinados ás condições da extensão. Por isso a sua acção estará sujeita ás leis physico-mathematicas, relativas ao tempo e ao espaço. Como tambem, se o sujeito da alteração está afastado e o \emph{meio} não se pode alterar num instante, é claro que as qualidades sensiveis não podem actuar no sujeito sem o concurso de movimentos locaes, e, d'este modo, os seus effeitos podem ser apreciados tambem pela \emph{Mechanica}. Mas é necessario sempre ter em vista esta grande verdade: embora as acções das qualidades sensiveis possam ser acompanhadas por movimentos locaes, \emph{não consistem} comtudo nesses movimentos. --- Advertimos finalmente que, pela analogia entre as formas \emph{substanciaes} e \emph{accidentaes}, devemos reconhecer nessas qualidades uma \emph{subordinação} ou \emph{continencia virtual}, isto é, que uma qualidade superior contém \emph{virtualmente} a inferior, e a excede. Por isso, ainda que a luz e a electricidade sejam \emph{distinctas} do calor, todavia o mesmo raio, a mesma onda pode ser, ao mesmo tempo, luminosa, calorifica e dotada de energia chimica, como se vê nos espectros da luz solar. Esta multiplicidade de energias, que a mesma onda, mechanicamente simples, pode encerrar, deriva do \emph{principio formal}, mais ou menos elevado, conforme a natureza do sujeito.}.

%%>>> t148

\pg{232} \pg{233} \tese{148}{Divisão do objecto da percepção sensitiva.}
O objecto da \emph{percepção sensitiva} é \emph{proprio}, \emph{commum} e \emph{concomitante}.

a) O objecto \emph{proprio} é o que póde ser percebido por um só sentido, e é constituido pelas \emph{qualidades sensiveis}, ou, em geral, por algumas propriedades singulares, cada uma das quaes se refere a um só sentido, externo ou interno \footnote{Cada sentido, tanto \emph{externo} como \emph{interno}, tem o seu objecto \emph{proprio}, como indicámos na Secção II d'este tratado.}.

b) O objecto \emph{commum} é o que é percebido por mais de um sentido, e é formado pelas qualidades \emph{materiaes}, que existem na \emph{quantidade} e se referem a varios sentidos. As principaes são: a \emph{extensão}, a \emph{figura}, o \emph{movimento}, a \emph{quietação} e o \emph{numero} \footnote{E' claro que um objecto \emph{commum} se refere a varios sentidos, mas sob diversos aspectos. Cada um d'estes diversos aspectos representa o objecto \emph{proprio} de um dos sentidos relativos. Por isso o objecto \emph{commum} percebe-se, mas só emquanto é percebido o \emph{proprio}. --- A \emph{grandeza} e a \emph{figura} são attingidas pela \emph{vista} e pelo \emph{tacto}. O \emph{numero} é attingido pela \emph{vista} e pelo \emph{sentido commum}, que adverte a successão dos movimentos. O \emph{movimento}, como tambem a \emph{quietação}, é objecto \emph{commum} dos sentidos da \emph{vista} e do \emph{tacto}. --- A estes cinco sensiveis \emph{communs} reduzem-se o \emph{tempo}, o \emph{estado}, a \emph{unidade}, a \emph{distancia} e a \emph{vizinhança}. Todos estes sensiveis \emph{communs} se reduzem á \emph{quantidade}. --- E, como as \emph{qualidades sensiveis}, que constituem o objecto \emph{proprio}, se encontram na \emph{quantidade}, ou \emph{extensão}, podemos dizer que os objectos \emph{proprios} e os \emph{communs} existem na \emph{quantidade}.\par Desenvolvamos este ponto. --- As \emph{qualidades sensiveis}, --- o \emph{calor}, o \emph{sabor}, o \emph{som}, etc., --- constituem o objecto \emph{proprio} dos sentidos, possuem a efficacia de actuar nos mesmos sentidos. Sendo, porém, formas \emph{accidentaes} e \emph{materias}, não podem essas qualidades existir senão na materia e, propriamente, na \emph{quantidade continua}, como no seu sujeito \emph{immediato}, e portanto não podem operar nem actuar nos sentidos, senão conformarem a sua acção com as condições da quantidade, e por isso a sua acção será \emph{subjectivamente} extensa. Diz S. Thomaz: « Qualitates sensibiles movent sensum \emph{corporaliter} et \emph{situaliter} » (In II \emph{De Anima}, l. 13). --- Se, porém, é certo que não ha um sentido que se refira \emph{primariamente} á quantidade (porque, se o houvesse, não teria especificação), é tambem certo que \emph{todos} os sentidos se referem a ella \emph{secundariamente}, isto é, por meio das \emph{qualidades sensiveis}, que são o objecto \emph{proprio} de cada um (porque, como dissemos, essas \emph{qualidades} só podem existir na \emph{quantidade continua}). Portanto a \emph{quantidade continua} é sensivel \emph{commum}, porque é o \emph{substrato} de todos os sensiveis \emph{proprios}, e pode ser percebida, mais ou menos distinctamente, por todos os sentidos. Ouçamos o S. Doutor: « \emph{Qualitates}, quae sunt \emph{propria} objecta sensuum, sunt \emph{formae} in \emph{continuo}; et ideo oportet quod ipsum continuum, in quantum est \emph{subjectum} talibus qualitatibus, moveat sensum \emph{non per accidens}, sed sicut \emph{per se subjectum} et \emph{commune} omnium sensibilium qualitatum » (\emph{De sensu et sensato}, l. 2). --- A' \emph{quantidade continua} referem-se e reduzem-se tambem, de algum modo, os \emph{sensiveis communs}. Na verdade, a \emph{grandeza} e o \emph{numero} são apenas especies da \emph{quantidade}; a \emph{figura}, por ser o termo da grandeza, é uma \emph{qualidade} da \emph{quantidade}; o \emph{movimento} importa a \emph{distancia} e a \emph{approximação} de um ser com relação a outro (\emph{Ibid.}, \emph{cf.} \emph{Sum. Theol.}, p. I, q. 78, a. 3 ad 2).}.

c) O objecto \emph{concomitante} é o que não se attinge em si, mas num outro objecto, a que está unido. Tal é a \emph{substancia}; por quanto, por meio dos sentidos percebemos as qualidades ou propriedades corporeas, não \emph{abstractas} mas \emph{concretas}, e as qualidades são \emph{concretas}, emquanto adherem a um sujeito, que é a \emph{substancia} \footnote{O objecto \emph{concomitante}, que os escholasticos chamam \emph{sensivel per}\pgn{235}\emph{accidens}, não actua por si no sentido, de que é sensivel \emph{per accidens}, não é percebido \emph{per se}, mas encontra-se no \emph{sensivel per se} e pode ser percebido por um outro sentido ou por uma outra faculdade perceptiva, sensitiva ou intellectual, do sujeito sensiente. Assim o \emph{sabor} de uma maçã \emph{vermelha} é sensivel \emph{per accidens} da \emph{vista}, e a \emph{substancia} d'ella é sensivel \emph{per accidens} de todos os sentidos.\par Como se vê, a divisão do objecto da percepção \emph{sensitiva} é muito racional. --- \emph{Objecto} da percepção sensitiva é tudo o que pode ser percebido pelos nossos sentidos, e pode ser percebido pelos nossos sentidos tudo o que pode actuar nelles. Ora um ser pode actuar nos sentidos --- ou \emph{por si mesmo}, --- ou por \emph{meio} de \emph{outro} ser, a que está unido e que actua por si mesmo. Se actua \emph{por si mesmo}, diz-se objecto \emph{per se}, e divide-se em \emph{proprio} e \emph{commum}. Se actua por \emph{meio de outro}, diz-se objecto \emph{concomitante} ou \emph{per accidens}.\par S. Thomaz expõe muito claramente as \emph{tres} especies do objecto da percepção \emph{externa} com as seguintes palavras: « Similitudo alicujus rei est in sensu \emph{tripliciter}. Uno modo \emph{primo} et \emph{per se}, sicut in visu est similitudo colorum, et aliorum sensibilium \emph{propriorum}. Alio modo \emph{per se}, sed non \emph{primo}, sicut in visu est similitudo \emph{figurae} vel \emph{magnitudinis}, et aliorum sensibilium \emph{communium}. Tertio modo, nec \emph{primo}, nec \emph{per se}, sed \emph{per accidens}, sicut in visu est similitudo hominis, non in quantum est homo, sed in quantum huic colorato \emph{accidit} esse hominem » (\emph{Sum. Theol.}, p. I, q. 17, a. 2). --- Noutro logar (\emph{Sum. Theol.}, p. I, q. 78, a. 3 ad 2) o S. Doutor explica as palavras, que acabamos de transcrever. O objecto \emph{proprio} é o que, \emph{por si mesmo} e \emph{antes} ou \emph{independentemente} dos outros objectos, actua na faculdade sensitiva e a determina ao acto; tal é a côr com relação aos olhos. O objecto \emph{commum} é o que \emph{modifica} a acção da \emph{quantidade} na faculdade sensitiva, e por isso actua \emph{por si mesmo} na faculdade, mas \emph{conjunctamente} com a quantidade e \emph{dependentemente} da quantidade; tal é a \emph{figura} de uma superficie, a qual actua \emph{por si mesma} nos olhos, \emph{modificando} a percepção da superficie (assim a percepção da superficie de uma certa figura differe da percepção da superficie de uma outra figura), mas não actua \emph{independentemente} da propria superficie. O objecto \emph{concomitante} é o que não actua na faculdade, nem \emph{por si mesmo}, nem \emph{dependentemente} de outro objecto, mas é apenas percebido emquanto se acha \emph{accidentalmente} unido a um sensivel \emph{proprio}; assim os olhos percebem o homem, não emquanto é homem, mas emquanto está como que envolvido na côr, que é o objecto \emph{proprio} da vista. --- Por isso o objecto \emph{commum} e o \emph{concomitante} não são percebidos senão \emph{dependentemente} do objecto \emph{proprio}; mas differem em que o \emph{commum} actua na facul\pgn{236}dade, modificando a impressão do objecto \emph{proprio}, ao passo que o \emph{concomitante} não actua de modo algum na mesma faculdade. --- A percepção do objecto \emph{proprio} e do \emph{commum} é \emph{distincta}; a do \emph{concomitante} é \emph{confusa}.\par A divisão do objecto dos sentidos, apresentada por Aristoteles, S. Thomaz e pelos escholasticos, não agradou a todos. Outros escriptores apresentaram outras. Limitamo-nos a citar as de \emph{Locke} e de \emph{Spencer}.\par \emph{a}) \emph{Locke} dividiu as \emph{qualidades sensiveis}, que constituem o objecto dos sentidos, em \emph{primarias} e \emph{secundarias}. As \emph{primarias} são as que se \emph{não podem separar} do proprio corpo, como a \emph{extensão}, a \emph{solidez}, a \emph{figura}, a \emph{mobilidade}; as \emph{secundarias} são as que se \emph{podem separar} do corpo, como a \emph{luz}, o \emph{calor}, a côr, o \emph{som}, o \emph{cheiro}, o \emph{sabor}, etc. --- As \emph{primarias} existem na realidade como os sentidos as percebem; as \emph{secundarias} não existem como os sentidos as percebem, e são combinações complexas da quantidade e do movimento.\par Esta divisão, como se vê, é \emph{incompleta} e \emph{arbitraria}. --- E' \emph{incompleta}, porque não falla do objecto \emph{concomitante}. --- E' \emph{arbitraria}, porque não se funda, como deveria, na essencial, primaria e immediata relação dos sentidos com o seu objecto \emph{proprio}; mas baséa-se numa erronea doutrina acerca da realidade do mesmo sensivel, e admitte como \emph{primario} o que é apenas \emph{secundario}, como é a relação dos sentidos com o sensivel \emph{commum}. A \emph{côr}, o \emph{sabor}, e, em geral, os sensiveis \emph{proprios}, são e devem chamar-se qualidades \emph{primarias}, porque \emph{determinam por si} e \emph{especificam} os sentidos relativos, e por isso são percebidas \emph{antes} das outras qualidades; ao passo que a \emph{figura}, a \emph{extensão}, etc., são e devem chamar-se \emph{secundarias}, porque não são percebidas senão \emph{por meio} das qualidades \emph{primarias}, e \emph{depois} d'ellas.\par \emph{b}) \emph{Spencer} dividiu as qualidades sensiveis em \emph{tres} classes, e são: as qualidades \emph{dynamicas}, as \emph{estatico-dynamicas}, as \emph{estaticas}. --- As qualidades \emph{dynamicas} são o \emph{som}, o \emph{cheiro}, etc., e, em geral, as que, para serem percebidas, não exigem \emph{actividade} da parte do sujeito sensiente (que, por isso, é totalmente \emph{passivo}). --- As \emph{estatico-dynamicas} são as \emph{resistencias}, que se encontram em certos objectos e suppõem \emph{actividade} e \emph{passividade} no sujeito sensiente (que, por isso, é \emph{activo} e \emph{passivo}, sob diversos aspectos). --- As \emph{estaticas} são a \emph{figura}, a \emph{situação}, e, em geral, as que o sujeito percebe de um modo simplesmente \emph{activo} (por isso o sujeito é só \emph{activo}).\par Esta divisão tambem é inacceitavel. Porquanto, além de ser \emph{incompleta}, porque não falla do objecto \emph{concomitante} ou \emph{per accidens}, é tambem \emph{arbitraria}, porque se funda --- não na relação estrictamente \emph{psychologica} das qualidades com os sentidos, que é primaria, \emph{immutavel}, natural e intrinseca, --- mas na relação meramente \emph{mechanica}, que é secundaria, variavel, accidental e extrinseca. --- Além disso, toda a qualidade sensivel, por se referir ao sentido, como o motor ao movido, como o agente ao paciente, é necessaria\pgn{237}riamente \emph{dynamica}. As \emph{estatico-dynamicas}, não podendo ser percebidas pelos sentidos, senão actuaram nelles, devem tambem dizer-se \emph{dynamicas}. Nenhuma qualidade sensivel pode ser \emph{estatica}, porque o sentido é essencialmente \emph{passivo}; a \emph{figura} e a \emph{situação}, que \emph{Spencer} chama qualidades \emph{estaticas}, devem tambem reduzir-se ás \emph{dynamicas}, visto não poderem ser percebidas senão emquanto se acham revestidas de alguma qualidade \emph{dynamica}, como é a \emph{luz}, a côr, o \emph{calor}, o \emph{cheiro}, a \emph{dureza}, etc.}.

149. \textbf{A percepção sensitiva é operação do composto.}

a) A \emph{percepção sensitiva} é propria do homem, emquanto é \emph{animal}; pois tal percepção é a operação \emph{especifica} do animal. Ora o \emph{animal} é o \emph{composto} de alma e de corpo. Logo a \emph{percepção sensitiva} é operação do \emph{composto}.

b) A \emph{percepção sensitiva}, considerada \emph{subjectivamente}, é, ao mesmo tempo, \emph{simples} e \emph{extensa}. Ora a \emph{simplicidade} e a \emph{extensão}, sendo propriedades \emph{contrarias}, --- não podem derivar --- da \emph{alma só}, porque esta é \emph{simples}, e do simples não pode derivar \pg{236}  o \emph{extenso}, --- nem do \emph{corpo só}, porque este é \emph{extenso}, e do extenso não pode derivar o \emph{simples}, --- mas devem derivar do \emph{composto}, porque o effeito deve ser proporcionado á causa. Logo a \emph{percepção sensitiva} é operação do \emph{composto}.

c) A \emph{percepção sensitiva}, considerada \emph{objectivamente}, attinge uma coisa material, dotada, ao mesmo tempo, de \emph{unidade} e de \emph{quantidade}. Ora uma coisa dotada de \emph{unidade} e de \emph{quantidade}, --- não pode ser attingida pela \emph{alma só}, que, sendo sim- \pg{237} , não pode apprehender a \emph{quantidade}, nem pelo \emph{corpo} só, que, sendo, \emph{quanto}, não pode apprehender a \emph{unidade}, --- mas deve ser attingida pelo \emph{composto}, porque só o \emph{composto}, por ser \emph{uno} e \emph{quanto}, pode apprehender um objecto \emph{uno} e \emph{quanto}. Logo a percepção \emph{sensitiva} é operação do \emph{composto} \footnote{O que dissemos tratando da \emph{simplicidade} da alma humana (\emph{Anthr.}, n. 30) e da \emph{materialidade} das faculdades sensitivas (\emph{Anthr.}, n. 73) dispensa o desenvolvimento dos argumentos, adduzidos para provar a nossa these. Diz S. Thomaz: « Sentire non est proprium animae, neque corporis, sed \emph{conjuncti}. Potentia ergo sensitiva est in \emph{conjuncto}, sicut in \emph{subjecto} » (\emph{Sum. Th.}, p. I, q. 77, a. 5). E noutro logar: « Quamvis animae sit aliqua operatio propria, in qua non communicat corpus, sicut intelligere; sunt tamen aliquae operationes \emph{communes} sibi et corpori, ut timere et irasci et \emph{sentire} et hujusmodi » (\emph{C. Gent.}, l. II, c. 57). Esta conclusão é o resultado logico e natural da analyse dos factos physiologicos da sensação. Esta analyse mostra que a \emph{unidade} e a \emph{extensão} pertencem realmente não só ao objecto sentido, mas tambem ao sujeito sensiente. Por isso o sujeito sensiente --- não pode ser exclusivamente a materia, porque esta, \emph{por si} só, não basta para sentir, --- nem exclusivamente a alma, porque o poder de sentir não pode existir fóra da materia, --- mas é o \emph{composto} de materia e de alma, isto é, o \emph{orgão animado}. Por isso diz S. Thomaz que o poder de sentir pertence á alma, não como a \emph{sujeito proximo} que por si só o exerça, mas como a \emph{principio ultimo} ou \emph{radical} d'onde deriva, porque é pela alma que o organismo vive e tem poder de sentir. Eis as suas palavras: « Omnes potentiae animae dicuntur esse animae non sicut \emph{subjecti}, sed sicut \emph{principii}; quia per animam habet \emph{conjunctum} quod tales operationes operari possit » (\emph{Sum. Theol.}, p. I, q. 77, a. 5). Tocámos neste ponto no tratado das \emph{faculdades}.\par \emph{Platão} e \emph{Descartes}, seguidos por A. \emph{Rosmini} e por alguns escriptores modernos, dizem que o corpo recebe a impressão dos objectos externos e que esta impressão é transmittida pelo proprio corpo á alma, a qual, por ser uma substancia simples e immaterial, é o verdadeiro e exclusivo sujeito da percepção sensitiva. D'ahi a celebre questão da \emph{ponte}, para explicarem o \emph{commercio} entre a alma e o corpo, --- questão que não teriam levantado; se tivessem considerado que a \emph{ponte} serve para unir duas coisas distantes, ao passo que as duas substancias, a alma e o corpo, estão por tal modo unidas que formam um só supposto e um só sujeito de operações.\par Examinemos algumas objecções contra a nossa these.\par a) Dizem: A percepção \emph{sensitiva} é uma operação exclusivamente propria da alma pelas seguintes razões: --- 1.º) porque a força sensitiva é \emph{simples};\pgn{238}--- 2.º) porque, aliás, dariamos logar ao materialismo; --- 3.º) porque, se a percepção \emph{sensitiva} fosse uma operação do corpo, a força sensitiva ficaria no proprio corpo, mesmo depois de separado da alma.\par A doutrina dos \emph{Cartesianos} e dos \emph{Rosminianos} é falsa. Por quanto, na opinião dos adversarios, a faculdade \emph{sensitiva} não seria \emph{material}, mas \emph{espiritual}; e assim entre ella e a intelligencia não haveria differença essencial. E, sendo assim, ou deveria negar-se aos animaes a energia sensitiva (como fez \emph{Descartes}), ou deveria attribuir-se á alma dos animaes a espiritualidade e a immortalidade (como fez \emph{Balmes}). Quem não vê o absurdo d'estas illações?\par Agora uma breve resposta aos argumentos dos adversarios.\par O \emph{primeiro} argumento é inefficaz para o scopo. Na verdade, quando os adversarios dizem que a energia sensitiva é \emph{simples}, póde entender-se --- ou que não é composta de partes \emph{essenciaes}, --- ou que não é \emph{extensa}. --- Na primeira hypothese, concedemos aos adversarios que a força sensitiva é \emph{simples}, mas isto não obsta a que ella resida num sujeito material; pois tambem a \emph{affinidade}, a \emph{gravidade}, e as outras qualidades dos corpos, são \emph{essencialmente} simples, e todavia existem em sujeitos materiaes. --- Na segunda hypothese, a força sensitiva não é \emph{extensa} \emph{em si}, mas é \emph{extensa} \emph{em razão do sujeito} corporeo, em que reside e de que depende.\par O \emph{segundo} argumento póde retorquir-se contra os adversarios. Porquanto, são elles que, baseando a espiritualidade da alma num fundamento falso, qual é a percepção sensitiva, dão occasião a que os materialistas neguem a mesma espiritualidade. --- A nossa alma é espiritual, não pelo facto de sentir (aliás tambem a alma do animal seria espiritual), mas pelo facto de produzir operações \emph{espirituaes}, como são as da intelligencia e da vontade.\par O \emph{terceiro} argumento é falso. Na verdade, a percepção \emph{sensitiva}, sendo operação \emph{vital}, depende por tal forma do principio da vida, que sem este não póde realisar-se. Ora o principio da vida é a alma.\par b) Insistem: Se a percepção \emph{sensitiva} fosse uma operação do \emph{composto}, não haveria identidade no sensiente; porque no composto uma parte é distincta da outra, e a operação de uma parte é extranha á da outra.\par Os adversarios suppõem que cada uma das partes do composto é o agente \emph{principal}. Este supposto é falso. O \emph{agente} principal é o \emph{todo}; as partes são apenas instrumentos. Ora o todo é \emph{uno}, e por isso o sensiente é dotado de identidade. --- Como se vê, os adversarios só reconhecem a \emph{unidade} e \emph{identidade}, proprias de um ser \emph{inextenso}.\par c) Alguns philosophos espiritualistas dizem: A percepção \emph{sensitiva} é essencialmente simples e indivisel, e por isso exige um sujeito \emph{simples}.\par Estes escriptores confundem a \emph{unidade} com a \emph{simplicidade}. \emph{A unidade} exclue a \emph{divisão}; a \emph{simplicidade} exclue a \emph{composição} ou \emph{extensão}. Tudo o que é \emph{simples} é \emph{uno}; mas nem tudo o que é \emph{uno} é \emph{simples}. --- A percepção \emph{sensitiva}, embora seja \emph{una}, é comtudo \emph{extensa} e \emph{divisivel}. --- E' \emph{extensa}. Por quanto, se pozermos a mão sobre uma meza de marmore, sentimos uma impressão de frio, e sentimol-a não num só ponto da mão, mas em toda a sua \emph{extensão}. --- E' \emph{divisel}. Com effeito, se da meza levantarmos um dedo, a percepção já não fica identica á primeira, mas é menos \emph{extensa}, porque\pgn{239}o dedo, que levantámos, já não sente a impressão do frio. A percepção sensitiva, pois, não é simples, mas extensa, como extenso é o organismo. --- Ora esta operação exige um sujeito proporcionado e por isso composto de um \emph{principio de unidade} e do um \emph{elemento de extensão} ou \emph{divisibilidade}; pois que, se não fosse composto de um \emph{principio de unidade}, o acto da percepção sensitiva não seria \emph{unico} e \emph{completo}, --- como tambem, se não fosse composto de um \emph{elemento de extensão} e \emph{divisibilidade}, a percepção sensitiva não seria, como effectivamente é, \emph{extensa} e \emph{divisivel}. Logo o sujeito da percepção \emph{sensitiva} não póde ser \emph{simples}.\par d) Os materialistas objectam: Corrompendo-se o orgão sensitivo, cessa ou fica viciada a respectiva percepção. Logo o sujeito da percepção sensitiva é só o orgão material, é só o corpo, e não o \emph{composto} de alma e de corpo.\par Resposta. Concedemos o \emph{antecedente}, e negamos o \emph{consequente} e a \emph{consequencia}. Por quanto, o facto de cessar ou ficar viciada a percepção sensitiva, pela corrupção do orgão respectivo, indica apenas que a energia sensitiva da alma precisa dos orgãos corporeos, como de instrumentos e de sujeitos, mas não significa, como os adversarios pretendem, que o orgão material, separado da força sensitiva, seja o sujeito da percepção. --- Todos os esforços, que os materialistas fizeram para reduzir os phenomenos psychicos á categoria dos factos meramente corporeos ou mechanicos, foram e serão sempre baldados. Affirma-o tambem um dos principes da sciencia experimental, \emph{Du Bois-Reymond}, por estas palavras: « Que nexo pode existir entre os movimentos de atomos determinados, que se realisam no meu cerebro, e certos factos especificos e innegaveis, como estes: eu \emph{sinto} uma dôr, \emph{experimento} um prazer, \emph{gosto} um sabor doce, \emph{aspiro} o perfume de uma rosa, \emph{ouço} o som do orgão, \emph{vejo} uma planicie verde? Não é possivel conceber de modo nenhum como a consciencia possa nascer do mero concurso dos atomos ». (\emph{Sobre os limites do conhecimento natural}, etc., Leipzig, 1884).}.

%%>>> t150

\tese{150}{A percepção sensitiva é verdadeiro conhecimento.}
\emph{Conhecimento} é uma operação \emph{immanente} pela qual o sujeito, depois \pg{234} \pg{235} \pg{238}  de ter recebido em si a imagem de um objecto, por meio d'esta apprehende o proprio objecto. Ora a \emph{percepção sensitiva} é uma operação \emph{immanente}, pela qual o homem, emquanto dotado \pg{239}  de vida sensitiva, depois de ter recebido a imagem \emph{sensivel} dos objectos materiaes, por ella apprehende esses objectos. Logo a percepção \emph{sensitiva} é verdadeiro \emph{conhecimento} \footnote{E' a doutrina de todos os escholasticos, sobretudo de S. Thomaz. Diz o S. Doutor: « Sentire \emph{cognoscere quoddam} est, et sentimus quandoque quidem in potentia, quandoque autem in actu » (\emph{In III De anima}, lect. 7). E noutro logar: « Species abstrahitur (\emph{derivatur}) a rebus ipsis: et per eam cognitio sensus continuatur ad ipsas res sensibiles » (\emph{De Verit.}, q. II, a. 6). --- E' necessario, para isso, não esquecer que todo e qualquer \emph{conhecimento} consiste numa certa \emph{semelhança} do objecto conhecido, realisada no cognoscente: « Omnis cognitio fit secundum similitudinem cogniti in cognoscente » (\emph{C. Gent.}, l. II, c. 77). Ora esta \emph{semelhança} não se realisa senão emquanto o cognoscente reproduz em si mesmo a coisa conhecida, por meio de uma imagem, que informa a faculdade cognitiva e representa a propria coisa. Diz S. Thomaz: « Dá-se uma assimilação entre duas coisas, emquanto ambas convêm na mesma natureza, mas não é esta a que é necessaria para o conhecimento. Dá-se outra assimilação entre duas coisas, emquanto uma é informada pela imagem da outra, e esta é necessaria para o conhecimento; assim a vista torna-se semelhante á côr, emquanto a especie ou imagem d'esta informa a pupilla: « \emph{Est quaedam assimilatio secundum convenientiam in natura... sed haec non requiritur ad scientiam. Est etiam quaedam assimilatio per informationem}\pgn{240}\emph{, quae requiritur ad cognitionem; sicut visus assimilatur colori, cujus specie informatur pupilla} » (\emph{I Dist.} 34, Q. 3, a. 1 ad 4). --- D'aqui se deduz a razão, pela qual os seres cognoscitivos são superiores aos que não conhecem, porque estes só possuem a propria forma \emph{physica} ou \emph{real}, pela qual são o que são e são collocados numa determinada especie, ao passo que os cognoscitivos possuem, além da propria forma \emph{physica} ou \emph{real}, outras formas \emph{intencionaes}, pelas quaes se tornam, não \emph{real} mas \emph{intencionalmente}, todas as coisas, que as formas representam. Diz S. Thomaz: « Cognoscentia a non cognoscentibus in hoc distinguuntur, quia non cognoscentia nihil habent nisi formam suam tantum, sed cognoscens natum est habere formam etiam rei alterius; nam species cogniti est in cognoscente » (\emph{Sum. Th.}, p. I, q. 14, a. 1).\par Os \emph{Rosminianos} objectam: A percepção \emph{sensitiva} não pode chamar-se \emph{conhecimento}. Por quanto, o \emph{conhecimento} exprime actividade, attinge o objecto externo, e é semelhante ao proprio objecto. Ora estas propriedades não competem á percepção \emph{sensitiva}; aliás esta distinguir-se-hia da \emph{intellectual} por uma differença, não \emph{qualitativa}, mas apenas \emph{quantitativa}.\par Resposta. As tres propriedades, que constituem a essencia do \emph{conhecimento} competem, com toda a justiça, á percepção \emph{sensitiva}. --- 1.º) A percepção \emph{sensitiva} exprime \emph{actividade}; porque a faculdade \emph{sensitiva}, embora seja \emph{passiva}, emquanto deve ser determinada pela acção do objecto, comtudo é dotada de energia vital, pela qual produz o seu acto de apprehensão. Diz S. Thomaz: « Sentire, quantum ad receptionem speciei sensibilis, nominat \emph{passionem}; sed quantum ad actum consequentem ipsum sensum, perfectum per speciem, nominat \emph{operationem} » (\emph{In I Sent.}, Dist. 40, q. 1, a. 1 ad 1). --- 2.º) A percepção \emph{sensitiva} attinge o objecto \emph{externo}; pois que, como veremos, toda a percepção é \emph{objectiva}. --- 3.º) A percepção \emph{sensitiva} é semelhante ao proprio objecto percebido, emquanto a faculdade, informada pela imagem ou semelhança do objecto, conhece o objecto e conhece-o no proprio ser, que o objecto possue fóra da faculdade. Diz S. Thomaz: « Similitudo rerum est in sensu » (\emph{Sum. Theol.}, p. I, q. 17, a. 2). E noutro logar: « Oculus per speciem lapidis, quam habet in se, cognoscit lapidem secundum esse, quod habet extra oculum » (\emph{Ib.}, q. 14, a. 6). --- Competindo á percepção \emph{sensitiva} as propriedades, que constituem a essencia de \emph{conhecimento}, é claro que a essa percepção deve convir tambem o nome de \emph{conhecimento}. --- Todavia, do facto de a percepção \emph{sensitiva} possuir os requisitos, que constituem a essencia do conhecimento, não se segue que ella se distinga da percepção \emph{intellectual} por uma differença apenas \emph{quantitativa}, e não \emph{qualitativa}. Se as duas percepções convêm no \emph{genero} de \emph{conhecimento}, differem na \emph{especie}, e por isso na \emph{qualidade}. Porque o homem convém com o bruto no \emph{genero} de \emph{animal}, poderá dizer-se com razão que a differença entre um e outro é \emph{exclusivamente quantitativa}? Não, por certo. --- Como se vê, estes escriptores são dignos discipulos do mestre, que (\emph{Nuovo Saggio} ec., vol. I, p. 34 e seg.) reserva exclusivamente para a \emph{intelligencia} o \emph{conhecimento}, e admitte que o animal \emph{sente}, mas \emph{não conhece}. E' uma affirmação arbitraria e contraria ao senso commum e á linguagem universal. Todos dizem que o pintainho \emph{conhece} a mãe, que o cão \emph{conhece} o dono, etc.}.

%%>>> t151

\tese{151}{Divisão da percepção sensitiva.}
A percepção \emph{sensitiva} divide-se em \emph{externa} e \emph{interna}, --- A \emph{externa} é o acto proprio dos sentidos \emph{externos}, pelo qual o homem conhece de um modo \pg{240} \pg{241}  \emph{material} os entes \emph{materiaes} e \emph{externos}. --- A \emph{interna} é o acto dos sentidos \emph{internos}, pelo qual o homem conhece de um modo \emph{material} os seus actos e estados \emph{sensitivos} e \emph{internos} \footnote{Muitos escriptores modernos negam uma distincção \emph{real} entre a percepção sensitiva \emph{externa} e a \emph{interna}, e afirmam que toda e qualquer percepção só representa os varios estados da nossa consciencia.\par E' manifesto o erro d'esses escriptores. A distincção entre a percepção sensitiva \emph{externa} e a \emph{interna} não é arbitraria, mas funda-se numa rigorosa observação dos \emph{orgãos} e dos \emph{factos}. --- Na verdade, a analyse \emph{anatomica} e \emph{physiologica} mostra que os \emph{orgãos} periphericos dos sentidos \emph{externos} differem completamente, na estructura e na situação, dos \emph{orgãos} dos sentidos \emph{internos}; e uma completa differença nos orgãos importa uma perfeita distincção nas respectivas funcções. --- Emquanto aos \emph{factos}, attesta-nos a consciencia que existe uma differença radical entre a percepção d'um acto \emph{interno} ou d'um estado affectivo e a percepção de um objecto \emph{externo}. Quando experimentamos uma dôr, percebemos (se não se der o caso da allucinação) que ella affecta o nosso organismo, que somos nós que soffremos, e não é um agente externo, que a pode ter causado. Como, pelo contrario, quando percebemos a figura, a extensão, etc., de um corpo externo, que actua nos nossos sentidos, não attribuimos essa figura, ou extensão, a nós mesmos, mas a uma coisa que está fóra de nós. Quando a ponta de uma agulha se nos espeta num braço, percebemos claramente que a ponta pertence á agulha e a dôr ao nosso braço. E quando tocamos com a extremidade do dedo uma pequena bola de marfim, attribuimos á impressão do dedo a forma concava, e á bola a forma convexa. --- Portanto, as percepções \emph{externas} e as \emph{internas} são duas series de factos inconfundiveis e irreductiveis ao mesmo principio proximo. Confundir, identificar todos esses factos, e reduzil-os a meros estados da consciencia, a meras modificações do ser sensitivo, com a exclusão de toda e qualquer \emph{representação} dos objectos \emph{externos}, é uma verdadeira arbitrariedade, reprovada pela razão e pela experiencia.}.

%%>>> art-22

\section*{Artigo II.\\ Percepção sensitiva externa.}
\addcontentsline{toc}{section}{Artigo II. --- Percepção sensitiva externa.}

%%>>> t152

\tese{152}{Percepção sensitiva externa.}
Percepção sensitiva \emph{externa} é, como dissemos, o acto proprio dos sentidos \emph{externos}, pelo qual o homem conhece de um modo \emph{material} ou \emph{concreto} os entes \emph{materiaes} e \emph{externos} \footnote{O assumpto, de que trata este artigo, é de capital importancia, porque é a base de todo o edificio intellectual. Como devemos restringir as nossas considerações entre os limites de um Compendio, remettemos os estudantes e leitores ás classicas obras do P. Salis Seewis (\emph{Della Conoscenza Sensitiva}) e de M. Farges (\emph{L'objectivité de la perception des sens externes}).}.

%%>>> t153

\tese{153}{Objecto da percepção sensitiva externa.}
O \emph{objecto} da percepção sensitiva \emph{externa} é constituido pelos entes \emph{sensiveis} e \emph{externos}. --- Estes objectos, como dissemos, actuam nos sentidos \pg{242}  tidos \emph{externos} pelas \emph{propriedades sensiveis}, de que são dotados e que são o objecto \emph{directo} dos mesmos sentidos \footnote{O \emph{objecto} da percepção sensitiva \emph{externa} é constituido pelos entes \emph{materiaes} ou \emph{sensiveis} e \emph{externos}. Diz S. Thomaz: « Est... sensus quaedam potentia \emph{passiva}, quae nata est immutari ab \emph{exteriori sensibili} » (\emph{Sum}. \emph{Th}., p. I, q. 78, a. 3). E, como os sentidos externos não se determinam ao acto senão pela acção do objecto, segue-se d'ahi que este deve estar \emph{presente}. Diz o mesmo S. Doutor: « Sensus secundum suam propriam rationem non est cognoscitivus nisi \emph{praesentium} » (\emph{De sensu et sensato}, l. 1). Por isso, se algumas faculdades \emph{sensitivas internas} (como a \emph{imaginação}, a \emph{memoria}) percebem ou reproduzem objectos \emph{ausentes}, esta virtude não a possuem ellas emquanto são sensitivas, mas emquanto participam, por uma analogia, da faculdade \emph{intellectual}. Continua S. Thomaz: « Quod autem sit aliqua virtus sensitivae partis se extendens ad alia \emph{quae non sunt praesentia}, hoc est secundum similitudinariam participationem rationis vel intellectus » (\emph{l. c.}).}.

%%>>> t154

\tese{154}{A percepção sensitiva externa exige a união do objecto com a faculdade.}
A percepção \emph{sensitiva}, como demonstrámos, é um verdadeiro \emph{conhecimento}. Ora todo o \emph{conhecimento}, seja qual for a sua especie, exige a união do objecto percebido com a faculdade perceptiva. Logo a percepção \emph{sensitiva} exige a união do objecto com a sua faculdade \footnote{A percepção sensitiva externa, como toda e qualquer percepção, suppõe e exige necessariamente tres elementos: um \emph{sujeito}, capaz de perceber, um \emph{objecto}, capaz de ser percebido, e a \emph{presença} do objecto no sujeito. --- O \emph{sujeito}, como vimos, não é só a alma, nem só o corpo, mas o \emph{orgão animado}. --- O \emph{objecto} não é todo e qualquer ente material, mas só o que, além da quantidade, é dotado de força sufficiente para actuar nos sentidos. --- A \emph{presença} do objecto no sujeito é indispensavel, porque, como dissemos, toda a percepção faz-se emquanto o ente conhecido se encontra na faculdade cognoscitiva. Diz S. Thomaz: « Cognitio non dicit effluxum a cognoscente in cognitum, sicut est in actionibus naturalibus, sed magis dicit existentiam cogniti in cognoscente ». (\emph{De Verit}., q. 2, a. 5 ad 15).}.

%%>>> t155

\tese{155}{O objecto deve unir-se com a faculdade por uma sua imagem.}
A percepção \emph{sensitiva} exige a união do objecto com a propria faculdade. O objecto deve unir-se com a faculdade --- ou na sua realidade material, ou por meio de uma sua imagem. Mas não póde unir-se com a faculdade na sua realidade material; pois que os corpos são naturalmente impenetraveis. Logo deve unir-se-lhe por meio de uma sua imagem. --- Esta imagem, indispensavel para a percepção, é formada pela acção que o proprio objecto exerce na faculdade perceptiva, é material e concreta, e chama-se \emph{especie impressa} \footnote{Esta \emph{imagem}, ou \emph{especie}, ou \emph{representação}, foi chamada \emph{impressa}, porque é o proprio objecto, que a \emph{imprime} na faculdade. --- Fallemos, pois, da \emph{necessidade} da especie \emph{impressa}, como tambem da sua \emph{origem} e \emph{natureza}.\par \emph{a}) \textsc{NECESSIDADE} DA ESPECIE IMPRESSA. --- Sendo a percepção \emph{sensitiva} um\pgn{243}verdadeiro \emph{conhecimento}, devemos conceder-lhe tudo o que é indispensavel para a essencia \emph{generica} do \emph{conhecimento}. Ora, para essa essencia é indispensavel a união do objecto com a faculdade; porque o conhecimento, como dissemos, consiste propriamente nisto --- que o objecto conhecido se encontra na faculdade cognoscitiva. E se o objecto não pode penetrar na faculdade com a sua substancia material, pode e deve penetrar nella com uma sua \emph{imagem}, que o represente. Esta imagem completa a faculdade em ordem ao acto, e por isso se torna \emph{comprincipio} da propria percepção. Diz S. Thomaz: « Nós sentimos ou entendemos actualmente alguma coisa, emquanto o nosso sentido ou a nossa intelligencia é actualmente informada pela imagem da propria coisa sensivel ou intelligivel » (\emph{Sum}. \emph{Theol}., p. I, q. 14, a. 2). E noutro logar: « Para a visão, quer sensivel quer intellectual, são indispensaveis duas coisas: a energia visual, e a união da coisa vista com a respectiva faculdade. Com effeito, não pode haver visão, senão emquanto a coisa vista se encontra na faculdade. E isto percebe-se facilmente, quando se trata de coisas corporaes, que não podem encontrar-se na faculdade com a propria substancia, mas só com uma sua imagem; assim nos nossos olhos não se encontra a substancia da pedra, mas só a sua imagem, pela qual se vê actualmente a pedra » (\emph{Sum}. \emph{Theol}., p. I, q. 12, a. 2).\par Objectam: As especies \emph{impressas} não são necessarias para a percepção \emph{sensitiva}, porque, se o fossem, a percepção seria meramente \emph{subjectiva}.\par Resposta. A percepção \emph{sensitiva} seria meramente \emph{subjectiva}, se a especie \emph{impressa} fosse o \emph{termo} da propria percepção. Mas não é assim. A especie \emph{impressa}, como veremos, não é o \emph{termo}, mas é o \emph{principio} da percepção; isto é, o que se percebe directamente não é a especie \emph{impressa}, mas é o \emph{objecto}, que a especie representa. Sendo o \emph{principio} da percepção, a especie é necessaria, e a propria percepção, se é \emph{subjectiva}, emquanto é acto \emph{immanente} da faculdade, é \emph{objectiva}, emquanto ao objecto que se percebe.\par Se não se admitte a especie \emph{impressa}, é impossivel explicar a verdadeira percepção \emph{externa}. Os sentidos \emph{externos} são por si indifferentes para a percepção d'este ou d'aquelle objecto sensivel. Se, portanto, percebem um objecto, e não um outro, a razão é que foram determinados pela especie ou imagem d'aquelle objecto, e não de outro. --- Rejeitada a especie \emph{impressa}, abre-se a porta ao \emph{subjectivismo} de Kant, ou ao \emph{idealismo} de Schelling, ou ao \emph{pantheismo} de Hegel, ou ao mais completo \emph{scepticismo}.\par b) \textsc{ORIGEM} DA ESPECIE IMPRESSA. --- A \emph{imagem} deriva da acção do \emph{objecto} no sujeito ou na faculdade cognoscitiva. Por quanto, a união do objecto com o sujeito não é uma \emph{juxtaposição} de dois entes, mas é a união do \emph{agente} com o \emph{paciente}. Os corpos, se não podem penetrar na faculdade perceptiva pela sua substancia material, penetram nella pela sua \emph{operação}, porque \emph{o termo da operação do agente encontra-se no paciente}. Penetrando na faculdade, a\pgn{244}operação de um agente sensivel, que é \emph{extensa}, \emph{luminosa}, \emph{figurada}, etc., manifesta alguma coisa do proprio agente; pois a operação é uma \emph{semelhança}, embora parcial ou inadequada, do seu principio. D'este modo o objecto, pela sua operação, \emph{imprime} e reproduz na faculdade a propria \emph{imagem}, ou \emph{especie}, ou \emph{figura}, ou \emph{semelhança}, a qual, unida com a faculdade, constitue um só ser na ordem \emph{intencional}, um só e adequado principio de operação, composto de \emph{potencia} e de \emph{acto} (\emph{Sum}. \emph{Theol}., p. I, q. 55, a. 1 ad 2). --- E nada mais natural. A operação do agente, informando a faculdade perceptiva, deve produzir e produz nesta uma \emph{paixão}, isto é, uma \emph{impressão}, que o representa, como o sinete imprime e reproduz na cera, em que actua, as armas, nelle gravadas; porque, como diz S. Thomaz, a paixão é sempre semelhante á acção: « \emph{Agens agendo assimilat sibi patiens} » (\emph{De Anima}, II, l. 10). Portanto, o primeiro effeito da acção do objecto no sujeito é o de tornar a faculdade \emph{semelhante} a si mesmo, não emquanto á \emph{substancia} ou \emph{natureza}, mas emquanto á \emph{imagem} ou \emph{representação}. --- A razão, por que os sentidos recebem a \emph{forma} ou \emph{imagem} do objecto material, e não a sua \emph{quantidade}, é exposta por S. Thomaz nas seguintes palavras: « O sentido recebe a \emph{especie} sem a materia, como a cera recebe a impressão das armas, gravadas no sinete, sem receber o ouro ou o ferro, de que o sinete é composto. E isto verifica-se em todo o paciente. Por quanto, o paciente recebe do agente alguma coisa, que lhe pertence, ou a elle se refere \emph{emquanto é agente}, e não emquanto possue outras propriedades. Ora o agente é tal, isto é, opéra em virtude da sua \emph{forma}, e não da sua \emph{materia}, se a tem. Por isso, todo o paciente deve receber e recebe a \emph{forma} do agente, sem receber a sua \emph{materia} » (\emph{De Anima}, II, l. 24; cf. \emph{De Verit}., q. 10, a. 4). --- Esta semelhança, que o agente produz no paciente, é incontestavel não só quando se trata da \emph{vista} e dos objectos \emph{extensos}, mas tambem quando se trata dos outros sentidos e dos respectivos objectos. Diz S. Thomaz: « Simili modo patitur et uniuscujusque sensus ab eo quod habet \emph{colorem}, aut \emph{saporem}, aut \emph{sonum} ». (In II \emph{De Anima}, c. 12). Mas d'isto mais adiante. --- Só accrescentamos que, podendo os sentidos receber as formas de todos os objectos sensiveis, não é erro dizer que os sentidos se tornam todos os sensiveis (não de um modo \emph{physico}, mas \emph{intencional}). D'ahi o adagio: \emph{Sensus sentiendo fit omnia}. Diz S. Thomaz: « Anima sensitiva est quodammodo omnia sensibilia, sicut anima intellectiva est quodammodo omnia intelligibilia » (\emph{Sum}. \emph{Theol}., p. I, q. 14, a. 1).\par \emph{Malebranche}, \emph{Leibnitz}, e outros, admittem a especie \emph{sensivel}, mas dizem que ella é o resultado da energia interna da alma na occasião de um phenomeno sensivel, sem que este influa de modo algum na producção da mesma especie. --- E' claro que, se assim fosse, seria impossivel a percepção \emph{externa}. Por isso diz S. Thomaz que, se a minha mão fosse aquecida não pela acção do fogo, mas pela acção de Deus, eu sentiria o calor da mão, mas não o calor do fogo. Eis as suas palavras: « Si sensatio caloris in organo (manus)\pgn{245}ab alio agente fieret (sc. a \textsc{Deo}), tactus etsi sentiret calorem manus, non tamen sentiret calorem ignis » (\emph{De pot}., q. III, a. 7).\par c) \textsc{NATUREZA} DA ESPECIE IMPRESSA. --- Mas o que é esta imagem, indispensavel para a percepção e produzida pela acção do objecto no sujeito? --- A \emph{imagem} pode considerar-se --- ou \emph{objectivamente}, isto é, emquanto ao \emph{objecto}, que ella representa, --- ou \emph{entitativamente}, isto é, emquanto á \emph{entidade}, que a constitue. D'ahi duas respostas á pergunta.\par \emph{a}) Considerada \emph{objectivamente}, a \emph{imagem} ou especie \emph{impressa} é a \emph{semelhança} do objecto sentido, não na ordem \emph{real} ou \emph{physica}, mas na ordem \emph{intencional} ou \emph{do conhecimento}. Chamamos de novo a attenção sobre este ponto, que é capital. A \emph{especie} é e deve ser \emph{semelhante} ao \emph{objecto}, não emquanto a \emph{especie} deva possuir a mesma \emph{natureza}, que o objecto possue, mas emquanto a \emph{especie representa} o \emph{objecto}. Por isso, uma imagem \emph{espiritual} pode \emph{representar} uma coisa \emph{material}, ou \emph{concreta}. Diz S. Thomaz: « Entre a faculdade cognoscitiva e o objecto conhecido deve haver uma \emph{semelhança}, não emquanto á \emph{natureza} (de modo que a faculdade e o objecto \emph{convenham} na mesma natureza), mas só emquanto á \emph{representação}. Com effeito, a forma da pedra, que se encontra na alma, possue uma natureza muito diversa da natureza, que tem na materia, e todavia a forma da pedra, emquanto \emph{representa} a mesma pedra, leva a faculdade a percebel-a » (\emph{De Veritate}, q. 8, a. 11 ad 3). E mais adiante: « Não é necessario --- que o objecto conhecido possua a natureza do sujeito que conhece, --- ou que o seu modo de ser corresponda perfeitamente ao modo de ser, que a forma tem na faculdade perceptiva. Por isso não é impossivel que por meio de formas, que existem na mente de um modo \emph{immaterial}, se possam conhecer coisas \emph{materiaes} » (\emph{Ibid}., q. 10, a. 4). --- D'onde se vê o erro de \emph{Reid}, que sustenta não poder haver \emph{semelhança} entre o objecto externo e a especie sensivel. Se é verdade que entre estas duas coisas não pode haver e não ha semelhança \emph{univoca}, fundada na conveniencia ou participação da mesma natureza; é todavia falso que entre o objecto externo e a especie sensivel não possa haver e não haja uma semelhança \emph{analoga}, de \emph{representação}, resultado da impressão do objecto no sujeito. Ora, diz S. Thomaz, não é a \emph{primeira} semelhança, que se requer para a percepção, mas a \emph{segunda} (\emph{Dist}., 34, q. 3, a. 1 ad 4).\par \emph{b}) Considerada \emph{entitativamente}, a imagem --- não é uma \emph{substancia}, embora subtilissima, que, destacando-se dos corpos, vóe pelo ar e vá ferir os sentidos, como sonharam \emph{Democrito} e \emph{Epicuro} nos tempos antigos, e \emph{Malebranche} nos modernos, --- mas é uma \emph{forma accidental}, que, unida á faculdade perceptiva, aperfeiçoa-a e completa-a intrinsecamente em ordem á operação. Por quanto, o \emph{acto} e a \emph{potencia} pertencem ao \emph{mesmo genero}, como explicámos. Se a faculdade perceptiva, que é a \emph{potencia}, é uma coisa \emph{accidental}, tambem a imagem impressa, que é o \emph{acto} d'essa potencia, não pode deixar de ser uma coisa \emph{accidental}. --- Mas esta forma \emph{accidental} é \emph{espiritual} ou \emph{material}? Para respondermos cabalmente a esta pergunta, observamos, com o Angelico Doutor, que a \emph{forma}, ou o \emph{ser}, de um objecto \emph{sensivel} pode considerar-se num triplice estado: em \emph{si mesma}, na \emph{intelligencia}, nos \emph{sentidos}. Considerada em \emph{si mesma}, no seu estado \emph{natural}, a \emph{parte rei}, a \emph{forma} é\pgn{246}inteiramente corporea, material e determinada pelas condições materiaes. Considerada na \emph{intelligencia}, a forma é inteiramente espiritual, immaterial e abstracta de todas as condições individuantes. Considerada nos \emph{sentidos}, a forma não é inteiramente corporea, porque abstrahe da materia ou da quantidade; mas não é inteiramente espiritual ou immaterial, porque não abstrahe das condições materiaes. Por isso, a forma, que existe nos sentidos, é \emph{intermedia} entre a que existe em \emph{si mesma} e a que existe na \emph{intelligencia}. Eis as palavras do S. Doutor: « Esse formae in \emph{imaginatione}, quod est quidem \emph{sine materia}, non autem \emph{sine conditionibus materialibus}, \emph{medium} est inter esse formae, quae est \emph{in materia}, et esse formae, quae est in \emph{intellectu} per abstractionem a materia et conditionibus materialibus » (\emph{Sum}. \emph{Theol}., p. I, q. 55, a. 2 ad 2). Portanto, a imagem ou especie \emph{sensivel}, --- por ser, ao mesmo tempo, cercada de condições materiaes, como o \emph{objecto material}, e abstracta da materia ou quantidade, como a \emph{especie intelligivel}, participa das qualidades d'esses dois extremos. E, como o ser \emph{intermedio}, comparado com um dos dois extremos, se parece com outro extremo, segue-se que a especie \emph{sensivel}, --- comparada com o objecto material, pode dizer-se \emph{espiritual}, porque não tem materia ou quantidade, como a tem o objecto, --- e, comparada com a \emph{especie intelligivel}, pode chamar-se \emph{material}, porque está cercada de condições materiaes, das quaes abstrahe a \emph{intelligivel}. --- Esta simples observação é sufficiente para explicar as palavras de S. Thomaz, que umas vezes chama \emph{espiritual} a forma sensivel, outras vezes \emph{material}. Em geral, os escholasticos chamam-na \emph{material}. O essencial é que se entenda bem que a especie \emph{sensivel} não tem materia ou quantidade, mas é \emph{cercada de condições materiaes}. Na verdade, devendo a entidade recebida num sujeito amoldar-se á indole do proprio sujeito, é claro --- que os sentidos, faculdades materiaes e corporeas, não podem receber a semelhança dos objectos percebidos senão de um modo material ou corporeo, --- e que a intelligencia, faculdade espiritual ou immaterial, não pode receber a semelhança dos objectos entendidos senão de um modo espiritual ou immaterial. Diz S. Thomaz: « Oportet quod sensus \emph{corporaliter} et \emph{materialiter} recipiat similitudinem rei quae sentitur, intellectus autem recipit similitudinem ejus quod intelligitur, \emph{incorporaliter} et \emph{immaterialiter} » (\emph{De Anima}, II, l. 12). Noutro logar, o S. Doutor explica-se, dizendo que os objectos são recebidos nos sentidos \emph{materialmente}, isto é, com \emph{as condições materiaes}, ainda que \emph{sem a materia}. Eis as suas palavras: « Sensus et imaginatio sunt vires organis affixae corporalibus; et ideo \emph{similitudines} rerum recipiuntur in eis \emph{materialiter}, id est \emph{cum conditionibus materialibus}, quamvis \emph{absque materia} » (\emph{De Verit}., q. 2, a. 5 ad 2).\par d) \textsc{OFFICIO} DA ESPECIE IMPRESSA. --- As opiniões dos philosophos ácerca do officio da especie \emph{impressa} são muitas e diversas.\par a) \emph{Malebranche}, \emph{Leibnitz}, e, em geral, os \emph{espiritualistas exagerados}, que affirmam ser a especie \emph{impressa} um producto exclusivo da energia da alma, concluem que a mesma especie é o verdadeiro \emph{objecto}, o verdadeiro \emph{termo} da percepção, é \emph{id quod cognoscitur}. Por isso a percepção é exclusivamente\pgn{247}\emph{subjectiva}, e é só por um acto de fé que julgamos a nossa percepção conforme com o mundo exterior.\par \emph{b}) \emph{Locke} admitte que a especie \emph{impressa} deriva da acção do objecto no sujeito, ainda que a sua theoria relativa á união da alma com o corpo não possa explicar logicamente a formação d'essa especie nas nossas faculdades pela acção de um phenomeno material. Emquanto ao \emph{officio} da especie, diz este escriptor que ella é o meio \emph{no qual} o objecto é percebido, é \emph{id in quo aliud cognoscitur}. Por isso a nossa alma conhece \emph{immediata} e \emph{directamente} a imagem \emph{subjectiva}, e d'esta, como de um effeito, eleva-se, pelo principio de \emph{causalidade}, ao objecto real, de que deriva, e que, portanto, é conhecido de um modo \emph{mediato} e \emph{indirecto}. --- A especie \emph{sensivel} é, para \emph{Locke}, como a imagem de uma pessoa, que se vê no espelho. O que se vê \emph{immediata} e \emph{directamente} é a imagem; e é nesta imagem que se vê a pessoa, a qual, por isso, não se vê em si mesma, não só vê de um modo \emph{directo} e \emph{immediato}, mas só de um modo \emph{indirecto} e \emph{mediato}. Em seguida, julgamos, pelo principio de \emph{causalidade}, que a pessoa deve ser tal, qual a imagem a representa.\par \emph{c}) Os \emph{escholasticos} ensinam que a especie \emph{sensivel} não é o \emph{termo} da percepção, nem é o meio \emph{em que} o objecto é percebido, mas é o meio \emph{pelo qual} a faculdade attinge o objecto, é \emph{id quo objectum cognoscitur}. Diz S. Thomaz: « \emph{Species quae est in visu non est id quod videtur, sed quo visus videt; quod autem videtur est color, qui est in corpore} » (\emph{De Anima}, III, l. 9). E noutro logar: « \emph{Species (impressa) non est illud quod sentimus, sed magis illud quo sensus sentit} » (\emph{Sum}. \emph{Th}., p. I, q. 85, a. 2). Portanto, a especie é um \emph{signal}, porque representa alguma coisa á faculdade perceptiva; mas é um signal \emph{formal}, e não \emph{objectivo}. O signal --- é \emph{formal}, quando leva a faculdade a perceber de um modo \emph{immediato} e \emph{directo} a coisa representada, sem que elle seja \emph{primeiramente} percebido; --- e é \emph{objectivo}, quando é percebido em si mesmo e depois leva a faculdade a perceber a coisa representada. --- Na doutrina dos \emph{escholasticos}, a faculdade não percebe a imagem sensivel, mas, estimulada e determinada por ella, reage e \emph{exprime} em si a mesma imagem, e, \emph{exprimindo} a imagem, attinge e percebe \emph{directamente} em si mesmo o respectivo phenomeno sensivel, ou o objecto material externo. Por isso a especie \emph{impressa} é uma causa instrumental, é uma excitação vital, um meio, ainda que indispensavel, para a faculdade produzir o seu acto de reacção e attingir o phenomeno material na sua realidade. --- Para a exacta comprehensão d'esta theoria, voltemos á comparação do espelho. Consideremos a nossa alma, não separada do espelho nem olhando para o mesmo, mas unida substancialmente a elle, animando-o e constituindo com elle um só sujeito sensiente. Supponhamos agora que uma imagem, derivada de um objecto material, venha imprimer-se na superficie d'este espelho, e que o espelho vivo, reagindo contra essa impressão, exprima a especie na direcção dos raios que emanam do objecto, e deveremos admittir que esse acto de reacção e de expressão vital attinge o proprio objecto externo.}.

%%>>> t156

\pg{243} \tese{156}{A imagem do objecto determina a faculdade á percepção.}

a) A faculdade sensitiva, sendo por si indifferente e indeterminada em relação aos innumeraveis objectos que estão ao seu alcance, para que possa perceber este ou aquelle objecto, deve sahir d'essa sua indifferença e indeterminação; pois que toda o agente opéra emquanto está em \emph{acto,} e não emquanto está em \emph{potencia}. Ora a faculdade \emph{sensitiva} sáe da sua indif- \pg{244}  ferença e indeterminação, quando o objecto externo imprime nella a sua \emph{imagem}; pois o ente, que existe em \emph{potencia}, é reduzido ao \emph{acto} pelo ente que já está em \emph{acto}, e todo o objecto, que existe na realidade, está em \emph{acto} relativamente á faculdade perceptiva. Logo a imagem do objecto determina a faculdade ao acto da percepção.

b) A experiencia confirma a these. Por quanto, se o objecto não estiver presente, de modo que possa operar na faculdade perceptiva, esta não produzirá o seu acto. Ora este phenomeno \pg{246}  seria inexplicavel, se o objecto, por meio de uma sua imagem, não determinasse a faculdade sensitiva á percepção. Logo a imagem do objecto determina a faculdade á percepção.

Esta determinação, proveniente da imagem do objecto, deve \pg{247}  affectar e affecta \emph{intrinsecamente} a faculdade; visto que uma indifferença \emph{intrinseca}, como é a da faculdade, só póde ser tirada por uma determinação \emph{intrinseca} \footnote{Toda a \emph{forma}, espiritual ou material, substancial ou accidental, determina o sujeito ou a uma certa especie, ou a uma certa operação; e esta especie ou operação é proporcionada ou semelhante á propria forma. Por\pgn{248}isso a imagem das coisas corporeas, embora material e accidental, determina a faculdade sensitiva, que a recebe, ao acto da percepção.\par D'onde se vê que a impressão, que o objecto produz no orgão corporeo, --- não é nem deve ser exclusivamente \emph{physica}, ou \emph{material}, como é a impressão do sinete na cera (aliás todos os corpos, mesmo os cadaveres, seriam capazes de sentir), --- mas deve ser, ao mesmo tempo, \emph{psychica}, emquanto o objecto, actuando no sentido, deve imprimir nelle a propria forma ou imagem, que é o principio \emph{formal} da percepção. --- A impressão \emph{physica} affecta o \emph{orgão}, embora destituido de energia sensitiva (assim um objecto externo imprime a propria imagem nos olhos de um animal morto); a \emph{psychica} affecta o \emph{sentido}, isto é, o \emph{orgão animado}. --- Todavia não devemos julgar que o objecto, actuando no orgão animado, produza nelle duas impressões \emph{justapostas}, --- uma \emph{physica}, outra \emph{psychica}, --- uma no \emph{orgão}, outra no \emph{sentido}. Não. O \emph{orgão} e o \emph{sentido} não são duas coisas \emph{justapostas}; o \emph{sentido}, isto é, a \emph{energia sensitiva} penetra por tal modo o \emph{orgão}, que constitue com elle um só principio sensiente, que é o \emph{orgão animado}. E, como ha um só principio sensiente, embora composto do orgão e da energia, tambem ha uma só impressão sensivel, proporcionada á faculdade perceptiva, e por isso nem exclusivamente physica, nem exclusivamente psychica, mas \emph{physico-psychica}.\par E', pois, manifesto o erro dos que affirmam que, para determinar a faculdade perceptiva ao acto, basta a \emph{mera} presença do objecto, embora este não influa na mesma faculdade. O objecto, que está presente mas que não influe na faculdade, é, emquanto ao effeito, como se não estivesse presente. Todavia é necessario que a faculdade saia da sua indifferença ou indeterminação, e não pode sahir della sem receber do objecto uma coisa que a determine. Essa coisa, que determina a faculdade, é a imagem sensivel.\par Esta é tambem a doutrina de S. Thomaz. Distingue o S. Doutor uma duplice \emph{mudança}, produzida pela acção do agente no paciente: uma \emph{natural} ou \emph{material}, outra \emph{espiritual} ou \emph{intencional}. A mudança é \emph{natural} ou \emph{material}, quando a forma do agente é recebida no paciente quanto ao \emph{ser natural} ou \emph{material}, assim a forma do calor é recebida no ferro aquecido ao lume. A mudança é \emph{espiritual} ou \emph{intencional}, quando a forma do agente é recebida no paciente quanto ao ser \emph{espiritual} ou \emph{intencional}; assim a forma da côr é recebida na pupilla dos olhos, que todavia não se torna colorida. Ora a impressão do objecto nos sentidos não produz apenas uma mudança \emph{natural} ou \emph{material} (aliás todos os corpos naturaes sentiriam, quando são alterados), mas produz tambem uma mudança \emph{espiritual} ou \emph{intencional}, pela qual a forma do agente é recebida no sentido. Eis as palavras de S. Thomaz: « Est \emph{duplex} immutatio: una \emph{naturalis} et alia \emph{spiritualis}. \emph{Naturalis} quidem, secundum quod forma immutantis recipitur in immutato secundum \emph{esse naturale}, sicut calor in calefacto; \emph{spiritualis} autem, secundum quod forma\pgn{249}immutantis recipitur in immutato secundum \emph{esse spirituale}, ut forma coloris in pupilla, quae non fit per hoc colorata. Ad operationem autem sensus requiritur immutatio \emph{spiritualis}, per quam intentio formae sensibilis fiat in organo sensus; alioquin, si sola immutatio \emph{naturalis} sufficeret ad sentiendum, omnia corpora naturalia sentirent, dum alterantur » (\emph{Sum}. \emph{Theol}. p. I, q. 78, a. 3). --- S. Thomaz chama \emph{natural} a mudança, ou impressão, que nós chamámos \emph{physica}, e chama \emph{espiritual} a \emph{psychica}.}.

%%>>> t157

\pg{245} \pg{248} \tese{157}{A faculdade, determinada pela imagem, percebe o objecto proprio externo, como é em si mesmo.}

a) Todos os homens, com pequenas excepções, estão convencidos de que os sensiveis \emph{proprios} existem nas substancias corporeas, e existem do mesmo modo por que são percebidos pela faculdade respectiva. Ora esse juizo do senso commum pode e deve ser attendido, quando não ha razões certas em \pg{249}  contrario, e, sobretudo, quando se trata de questões fundamentaes. Logo a faculdade sensitiva, determinada pela imagem, percebe o objecto \emph{proprio} externo, como é em si mesmo \footnote{Tratamos da \emph{objectividade} das nossas percepções \emph{sensitivas}. A questão refere-se principalmente á sensação \emph{externa}, mas interessa tambem a \emph{interna}; porque o \emph{sentido commum}, a \emph{imaginação} e a \emph{memoria} não percebem senão o que foi percebido, de algum modo, pelos sentidos externos. --- Os \emph{escholasticos} ensinaram sempre que na realidade, no mundo exterior, fóra do \emph{eu}, os objectos existem taes, quaes os referem os sentidos sãos e devidamente dispostos e applicados, e por isso disseram que não só os sensiveis \emph{communs}, mas tambem os sensiveis \emph{proprios}, --- as \emph{côres}, os \emph{sons}, etc. --- existem nos corpos externos como nós os percebemos (os sensiveis \emph{per accidens} não são percebidos \emph{em si} pelos sentidos). --- Veio \emph{Descartes}, e disse que nos corpos não existem as qualidades sensiveis, que nós costumamos reconhecer-lhes, e que as côres, os sons, e os outros sensiveis \emph{proprios} não passam de meros movimentos das particulas corporeas. --- \emph{Descartes} foi seguido por \emph{Locke}, que inventou a celebre divisão das qualidades em \emph{primarias} e \emph{secundarias}. --- As theorias de \emph{Descartes} e de \emph{Locke} agradaram a muitos escriptores, sobretudo aos naturalistas, e foram aproveitadas por alguns para fundamento do \emph{idealismo}. Assim, com o mesmo direito, com que \emph{Locke} tinha negado a realidade objectiva dos sensiveis \emph{proprios}, \emph{Berkeley} e outros negaram tambem a realidade das outras qualidades e de toda a entidade corporea, --- realidade, que \emph{Locke} tinha salvado. --- Neste \emph{subjectivismo} ou \emph{idealismo} ha muitos e diversos graus. --- As consequencias, que d'esses principios derivam, são: 1.º) a nossa intelligencia não pode conhecer de modo algum a natureza dos sensiveis \emph{proprios} só pelos dados fornecidos pelos sentidos; porque os sentidos não percebem essas qualidades, como são \emph{em si mesmas}, mas apenas como affectam os nossos orgãos; e por isso pertence á propria intelligencia interpretar o testemunho dos sentidos; --- 2.º) as nossas percepções \emph{sensitivas} não passam de \emph{illusões subjectivas}, ou de \emph{allucinações verdadeiras} (Taine); porque, ao passo que ás allucinações \emph{vulgares} nada corresponde fóra de nós, á percepção \emph{sensitiva} corresponde alguma coisa, que a intelligencia deve examinar e apreciar; --- 3.º) as nossas sensações não são \emph{imagens} ou \emph{semelhanças} dos objectos, mas são apenas \emph{signaes} (como as palavras são signaes das ideas), que não têm semelhança alguma com a coisa significada, e que devem ser interpretados ou explicados pela intelligencia; --- 4.º) o mais que podem admittir os que não seguem o \emph{idealismo} é que as qualidades sensiveis só existem nos corpos \emph{causaliter}, isto é, emquanto actuam nos sentidos, porque \emph{formaliter}, isto é, consideradas na sua natureza ou essencia, ellas não são taes senão pela sensação e na sensação actual. --- D'onde se vêm as differenças entre a doutrina antiga e a moderna. --- Os antigos admittiam que a especie sensivel não consiste no simples movimento, mas é uma qualidade produzida\pgn{250} pela acção do sensivel proprio e determinativa da propria faculdade. Os modernos affirmam --- que a mesma especie, quer se considere na origem de que deriva, quer no meio por que se propaga, quer no orgão em que é recebida, não passa de um \emph{movimento vibratorio}, --- e que o objecto externo não influe, de modo algum, no exercicio da actividade sensitiva. --- Os antigos ensinavam que os sensiveis \emph{proprios}, a luz com a côres, os sons, etc., têm em si mesmos a realidade, que os sentidos percebem; de modo que na faculdade sensitiva se encontra a imagem ou semelhança, não real mas intencional, do objecto percebido. Os modernos dizem que nada d'isso existe nos corpos, mas só existem vibrações e atomos diversamente figurados e dispostos, e por isso a tal representação ou imagem é um producto da intelligencia, ao qual não corresponde o objecto externo. --- Não temos uma razão plausivel para deixar os antigos e seguir os modernos. Tendo já explicado a natureza e o \emph{officio} da especie sensivel, resta-nos provar a \emph{objectividade} ou \emph{realidade} dos sensiveis \emph{proprios}. --- Lembramos que o objecto ou sensivel \emph{commum} se percebe, emquanto se percebe o objecto ou sensivel \emph{proprio}.\par Ao primeiro argumento, adduzido em favor da nossa these, respondem os adversarios que, em philosophia, não deve attender-se o juizo do \emph{senso commum}, que se tem muitas vezes enganado, como por ex. acerca do movimento da terra, da existencia dos antipodas, etc. --- Repetimos o que dissemos no texto. O juizo do \emph{senso commum} deve ser tido sempre em grande estima, emquanto não se provar \emph{com certeza} o contrario, sobretudo nesta questão capital, que é a \emph{objectividade} dos sensiveis \emph{proprios}, da cuja solução depende a vida ou a morte de todas as sciencias.}.

b) A faculdade sensitiva, como toda e qualquer faculdade perceptiva, é \emph{veraz}; o que se demonstra na \emph{Logica}. Ora a faculdade sensitiva não seria \emph{veraz}, se não percebesse os sensiveis \emph{proprios}, como são em si mesmos: o que é evidente. Logo a faculdade, determinada pela imagem, percebe o objecto \emph{proprio} externo, como é em si mesmo.

A faculdade percebe o objecto \emph{externo}, emquanto, reagindo, por um acto vital, sobre a imagem \emph{impressa}, a \emph{exprime}, --- isto é, emquanto, pelo acto e no acto da percepção, forma a imagem \emph{expressa} do proprio objecto \footnote{A conclusão é certa. A faculdade \emph{perceptivo-sensitiva externa} deve ter o seu \emph{objecto proprio}; porque toda a faculdade tem uma relação essencial e transcendental ao acto, e por isso ao objecto, que o acto attinge. Uma faculdade, que não tivesse um objecto \emph{proprio} e \emph{proporcionado}, seria uma contradicção. Os sentidos \emph{externos}, pois, devem ter um objecto \emph{proprio}, e este não pode ser senão o sensivel \emph{proprio}. Assim, a \emph{vista} tem por objecto \emph{proprio} a côr, e só em razão da côr pode perceber os sensiveis \emph{communs}, como a \emph{grandeza}, o \emph{movimento}, etc. Ora, se na realidade \emph{objectiva}, se fóra do \emph{eu}, não existissem os \emph{sensiveis proprios}, ou não existissem do mesmo modo por que os sentidos os referem, os sentidos não teriam o seu objecto \emph{proprio} e \emph{proporcionado}: o que é absurdo. --- E não só não teriam o objecto \emph{proprio}, mas nem seriam \emph{verazes}. Por quanto, elles referem que nos objectos \emph{externos} existem os sensiveis \emph{proprios}, a luz, o som, o calor, etc. Mas se esses\pgn{251} sensiveis \emph{proprios} não existissem nos objectos externos, ou não existissem do mesmo modo por que os sentidos os referem, os sentidos não seriam \emph{verazes}, porque a veracidade de uma faculdade consiste em referir uma coisa como é em si mesma. Diz S. Thomaz: « \emph{Veritas}... non sic est in sensu, ut sensus cognoscat veritatem, sed \emph{inquantum veram apprehensionem habet de sensibilibus}... Quod quidem contingit ex hoc, quod \emph{apprehendit}, vel \emph{judicat} res, \emph{ut sunt} » (\emph{Sum. Theol.}, p. I, q. 17, a. 2).\par Os adversarios sustentam que os nossos sentidos, embora sãos e devidamente dispostos, não percebem os sensiveis \emph{proprios} como são em si mesmos, e que por isso se enganam sempre com relação a esses sensiveis, e só se mostram verazes quando referem os sensiveis \emph{communs}.\par Resposta. --- Em primeiro logar, se, como dizem os adversarios, as côres, os sons, etc., não existissem na realidade, fóra do \emph{eu}; se o que se chama côr, som, etc., não passasse de velocissimos e minutissimos movimentos vibratorios, que escapam ao alcance dos sentidos, é claro que os sentidos não poderiam \emph{de modo algum} referir as côres, os sons, etc., por isso mesmo que de modo algum podem perceber esses movimentos, e assim não haveria \emph{illusão} ou \emph{engano} na percepção, mas simplesmente \emph{ausencia} ou \emph{negação} de percepção. --- Além d'isso: quando os adversarios affirmam que os sentidos referem os movimentos sob a forma \emph{ficticia} de côres, sabores, etc., visto não os poderem referir sob a verdadeira forma de movimentos, enunciam dois absurdos: --- que uma faculdade possa naturalmente alterar ou corromper o seu objecto, que ella não pode conhecer, revestindo-o de elementos ficticios, --- e que possa haver uma faculdade sem objecto \emph{proprio} e \emph{proporcionado}. --- Mais, dizem os adversarios que os sentidos se enganam com relação ao sensivel \emph{proprio}, mas não com relação ao sensivel \emph{commum}, porque o \emph{proprio} não é \emph{real}, como o é o \emph{commum}. Tambem esta affirmação é absurda. O sensivel \emph{commum} não é senão o \emph{proprio}, de alguma maneira modificado. O que é a \emph{grandeza}, emquanto percebida pela \emph{vista}, senão a côr mais ou menos extensa? Nenhum sentido, portanto, percebe o sensivel \emph{commum}, senão emquanto percebe o \emph{proprio}. O \emph{commum} suppõe o \emph{proprio}; aquelle segue o destino d'este. Por isso, se os \emph{communs} são reaes, tambem \emph{reaes} são os \emph{proprios}; e se os \emph{proprios} são \emph{ficticios}, tambem \emph{ficticios} são os \emph{communs}. Por isso são mais radicaes, mas mais logicos os que negam tambem a \emph{objectividade} do \emph{movimento}, que é um sensivel \emph{commum}. Diz \emph{Fouillée}: « O movimento: eis o grande idolo da sciencia moderna, mas não passa de um idolo. Pretende-se inutilmente que o adoremos, como o fundo mesmo da realidade. E' o Jupiter ou o Jehovah da physica! Mais racional e mais verdadeira é a opinião, que reduz, segundo o pensamento de Kant, d'Hamilton, e de Spencer, os movimentos externos, como as sensações internas, a meros symbolos de uma realidade escondida » (\emph{L'avenir de la Metaphysique}, p. 153). --- Mais ainda: se os sentidos se enganam acerca da \emph{objectividade} dos sensiveis \emph{proprios}, a intelligencia, que descobre e percebe o seu objecto nos elementos fornecidos pelos sentidos, não poderá conhecer coisa alguma com certeza, ou, melhor, estará sempre no engano e no erro, e assim nenhuma sciencia será possivel. --- Finalmente, ne analysarmos os argumentos, com que os\pgn{252}adversarios pretendem negar a \emph{veracidade} e a \emph{objectividade} dos nossos sentidos, veremos que elles não têm outra base a não ser o testemunho dos proprios sentidos. Logo, se o testemunho dos sentidos não é veraz para nós, quando affirmamos a \emph{objectividade} dos sensiveis \emph{proprios}, não pode ser veraz para os adversarios, quando contestam essa objectividade.\par Os escholasticos admittem a \emph{objectividade} de \emph{todos} os sensiveis \emph{proprios}, sem excepção, e por isso tambem das \emph{córes} e dos \emph{sons}. Não faltam, porém, escriptores que, embora admittam a nossa these fundamental, fazem comtudo as suas reservas com relação aos sensiveis \emph{proprios} da \emph{vista} e do \emph{ouvido}, dizendo que as \emph{córes} e os \emph{sons} têm apenas um estado \emph{psychico} e \emph{subjectivo} em nossos sentidos, e é por uma illusão natural que os referimos aos objectos externos. --- Poucas palavras para refutar esta opinião.\par E' certo que o \emph{senso commum} e a \emph{consciencia} attribuem aos objectos externos as \emph{córes} e os \emph{sons}. --- Será um \emph{absurdo} esta convicção? Assim dizem; mas dizem mal, porque um corpo colorido ou um sino sonoro não são conceitos incompativeis. --- Será \emph{obscuro} o conceito do som e das côres, como qualidades inherentes aos corpos? E' o que affirmam \emph{Descartes}, \emph{Locke} e muitos modernos; mas a obscuridade de um facto legitimamente averiguado, se é motivo para reconhecer a ignorancia ou a limitação da nossa intelligencia, não é motivo para negar a existencia do proprio facto. Se a obscuridade fosse uma razão sufficiente para negar uma verdade, poucas, muito pouco verdades ficariam de pé.\par Será contraria á \emph{sciencia} a nossa these da \emph{objectividade} das \emph{córes} e dos \emph{sons}? Assim o dizem, mas erram. Na verdade, a nossa these não é contraria a \emph{physiologia}, a qual se pode anatomisar todas as cellulas do cerebro, distinguir varias especies de nervos e perceber as vibrações d'estes nervos nas experiencias de vivisecção, não pode perceber na ponta do cinzel um phenomeno psychologico, e por isso não pode affirmar que a côr verde ou azul são meros estados da consciencia, phenomenos exclusivamente subjectivos. --- Mas, se a sciencia não averigua um phenomeno \emph{subjectivo}, pode averiguar um phenomeno \emph{objectivo}. E' certo. Pois bem, poderá a \emph{physica} dizer que os sons e as \emph{córes} não são realidades externas, porque nunca as poude averiguar? Tambem não. Por quanto, o physico serve-se dos seus apparelhos maravilhosos, mas d'esses apparelhos não se pode servir, sem se servir, ao mesmo tempo, dos cinco instrumentos primitivos, que nos deu Nosso Senhor para as observações externas e que são os sentidos corporaes. E se o physico se serviu da \emph{vista} para as \emph{córes}, e do \emph{ouvido} para os \emph{sons}, não nos poderá contradizer. Se, porém, não se serviu d'esses sentidos, a sua experimentação é incompleta e falsa. --- Mas, dizem, os sentidos da \emph{vista} e do \emph{ouvido} enganam-se. E como o sabem elles? Provaram, por ventura, de um modo cabal que as côres e os \emph{sons} não existem nos objectos, mas só nos sentidos? Não. E, se não merecem fé os sentidos da \emph{vista} e do \emph{ouvido}, porque haviam de merecel-a os outros? Ou todos falsos, ou todos verazes. --- Concluimos\pgn{253}que as sciencias experimentaes não podem de modo algum provar a \emph{subjectividade} dos \emph{sons} ed os \emph{córes}.\par Examinemos agora as principaes objecções dos adversarios.\par a) Dizem: O que é o \emph{som} e a cór? E' o que \emph{se} ouve e o que \emph{se} vê. Logo a \emph{essencia} do som e da côr consiste em \emph{ser ouvido} e em \emph{ser visto}; e por isso, onde não ha visão e não ha audição, não ha nem côr nem som.\par Resposta. Como é claro, os adversarios perguntam e respondem ao mesmo tempo, e da resposta, formulada a seu modo, tiram as consequencias que desejam. E' um methodo facil, mas pouco ou nada rasoavel. --- Feita esta observação, dizemos que a \emph{cór} não é só \emph{o que se} vé, mas é tambem \emph{o que pode ser visto}, como o \emph{som} não é apenas o \emph{que se ouve}, mas é tambem o que \emph{pode ser ouvido}. Logo a \emph{essencia da côr} e do \emph{som} consiste propriamente em ser \emph{visivel} e \emph{audivel}. Por isso, as \emph{córes} e os \emph{sons}, embora não sejam actualmente vistos ou ouvidos, não deixam de ser verdadeiras \emph{córes} e verdadeiros \emph{sons}, porque \emph{podem} sempre ser percebidos pela vista e pelo ouvido. Se assim não fosse, uma lampada accesa deixaria de ter luz, quando fecho os olhos, ou um barulho exterior acabaria num instante, quando e só porque tapo os ouvidos! Concluimos que as côres e os sons são o que são independentemente dos nossos sentidos, isto é, são \emph{objectivos} e \emph{reaes}. --- O erro dos adversarios nasce d'isto, que confundem o sentido e o objecto \emph{em acto} com o sentido e o objecto \emph{em potencia}. Se o sentido e o objecto se consideram \emph{em acto}, é claro que ambos se \emph{conservam} juntamente, ou juntamente \emph{perecem}, porque se um objecto sonoro não fôr \emph{actualmente soante}, não pode haver \emph{audição actual}, e vice-versa. Mas não se dá o mesmo, se o objecto e o sentido se consideram \emph{em potencia}, porque pode ser que se corrompa o sentido, sem que se corrompa o objecto, e vice-versa. E' o que S. Thomaz reprehendia nalguns naturalistas antigos, seguidos por alguns modernos, sobretudo por \emph{Bain}. Eis as suas palavras (In III, \emph{De An.}, l. 2): « Priores naturales non bene dicebant in hoc quia \emph{opinabantur nihil esse album aut nigrum nisi quando videtur}; \emph{neque saporem esse nisi quando gustatur}; et similiter de aliis sensibilibus et sensibus... Dicebant autem quodammodo recte et quodammodo non. Cum enim \emph{dupliciter} dicatur sensus et sensibile, scilicet \emph{secundum potentiam} et \emph{secundum actum}; de sensu et sensibili \emph{secundum actum accidit} quod ipsi dicebant: \emph{quod non est sensibile sine sensu}. \emph{Non autem hoc verum est} de sensu et sensibili \emph{secundum potentiam}. Sed ipsi loquebantur \emph{simpliciter}, idest, \emph{sine distinctione}, de his quae dicuntur \emph{multipliciter} ».\par b) Continuam: As operações de todos os corpos, simples ou compostos, reduzem-se a uma só: ao \emph{movimento}. As \emph{córes}, os \emph{sons}, o proprio \emph{calor}, são meros effeitos do movimento molecular. O \emph{som} é produzido por uma serie de vibrações moleculares, que se propagam atravez o ar ou os corpos ambientes; cada nota corresponde a um numero determinado de vibrações; cada tom dos instrumentos ou da voz humana corresponde a combinações mathematicas de vibrações, determinadas pelo calculo. A \emph{luz} e as \emph{córes} não passam de ondulações de um fluido subtilissimo, que provavelmente penetra todos os corpos e que se chama \emph{ether intermolecular}; e isto explica os phenomenos da \emph{interferencia}, os quaes consistem --- em que dois raios de luz homogenea e dotados da mesma velocidade, encontrando-se em sentido contrario\pgn{254}, se neutralisam e produzem obscuridade. O \emph{calor} é um movimento das mais pequenas partes do corpo; e as acções chimicas, por serem mechanicamente equivalentes á quantidade de calor produzido, transformado ou transportado, são outras especies de vibrações ou movimentos moleculares. Logo o \emph{som}, a \emph{luz}, como tambem o \emph{calor}, não são realidades \emph{objectivas}, isto é, existentes fóra de nós como nos apparecem, mas são meras affeições \emph{subjectivas}, e de \emph{objectivo} só têm o \emph{movimento}.\par Resposta. Podemos admittir todo o \emph{antecedente}, sem sermos obrigados a admittir o \emph{consequente}. Uma coisa é dizer que o \emph{som}, a \emph{luz}, o \emph{calor}, são \emph{acompanhados} de movimentos ou vibrações moleculares; outra coisa é dizer que a \emph{essencia} d'essas qualidades \emph{consiste} no mero movimento. E negamos o \emph{consequente}, porque \emph{nunca} esses escriptores poderão provar --- que nos phenomenos da \emph{luz}, do \emph{som} e do \emph{calor} não ha outra realidade \emph{objectiva} a não ser o movimento, --- e que as qualidades sensiveis, em geral, são \emph{formalmente} constituidas por um numero maior ou menor de vibrações, dotadas de maior ou menor velocidade.\par Mas procuremos uma resposta mais scientifica e mais cabal.\par \emph{Aristoteles}, que de muitos seculos precedeu os nossos modernos, disse tambem --- que todo o \emph{phenomeno} é o \emph{resultado} de um \emph{movimento}, --- que sem \emph{movimento} não ha phenomeno, --- que só pelo \emph{movimento} o phenomeno passa da potencia para o acto, --- que, supprimido o movimento, acaba immediatamente o estudo da natureza. E não só Aristoteles approvou antecedentemente as theorias modernas, mas, em muitos pontos, excedeu-as. A comparação, tida hoje por classica, das ondas sonoras com as ondulações circulares produzidas na agua por uma pedra, é invenção genial do grande Stagyrita (\emph{De Anima}, II, c. 8), e foi adoptada por S. Thomaz, que diz: « Sic autem contingit de immutatione aeris apud generationem soni, sicut de immutatione aquae, cum aliquid in aquam projicitur; manifestum est enim quod fiunt quaedam regyrationes in circuitu aquae percussae » (Ib., l. 16). Como tambem Aristoteles mediu, já no seu tempo, a \emph{intensidade} dos sons pelo numero mais ou menos rapido das vibrações (Arist., \emph{De Anima}, II, c. 8), e o mesmo ensinou S. Thomaz nestas palavras: « Sicut motus est causa soni, ita \emph{velocitas} motus est causa soni \emph{acuti}, et \emph{tarditas} motus est causa soni \emph{gravis} » (Id., l. 17). --- Mas onde se revelou, de um modo assombroso, o talento de Aristoteles foi na theoria da \emph{luz} e no systema das \emph{ondulações}, de que \emph{Huyghens} é tido como fundador e que hoje, depois dos trabalhos de \emph{Young} e de \emph{Fresnel}, é por todos recebido. Ensinou elle abertamente que a \emph{luz} não pode ser um corpo, nem emanação de um corpo, mas é um \emph{movimento} no meio especial, que elle chama \emph{diaphano}, e oppõe esta sua theoria á da \emph{emissão}, propugnada por \emph{Democrito} e pelos seus predecessores (\emph{De sensu}, c. III). --- Foi tambem Aristoteles quem estabeleceu analogias entre a \emph{luz} e o \emph{som}; falla da \emph{gamma} das \emph{córes}, e entrevê a possibilidade --- de \emph{exprimir as córes por numeros proporcionaes, como os accordos da musica} » (Ib., c. VII), --- e por isso de applicar as \emph{mathematicas} á sciencia da \emph{optica}, como \emph{Pithagoras} as tinha applicado á sciencia \emph{musical}. --- Ha mais. Fallando, no tratado de \emph{sensu} (c. VII), da \emph{luz}, do \emph{som}, do \emph{calor} e de todos os movimentos sensiveis, diz que podem \emph{contrariar-se} e até \emph{neutralisar-se}: « Si motus amborum\pgn{255}sint aequales, \emph{mutuo se obliterabunt}, cum ex eis unus non fiat », e assim formula o principio geral que explica os phenomenos da \emph{interferencia}, quando dois sons produzem silencio, e duas luzes obscuridade. --- Este systema das \emph{ondulações} é tão coherente com os principios relativos á natureza do movimento, que Aristoteles não só o applica á producção da \emph{luz}, do \emph{calor}, mas a todos os outros phenomenos sensiveis e á funcção de todos os nossos sentidos. Os proprios \emph{cheiros}, como a luz e as córes, são produzidos por vibrações, e não por emanações de corpusculos (\emph{De Anima}, II, c. VII). Tem razão \emph{Barthélemy St. Hilaire} de dizer: « A sciencia moderna não podia fallar melhor » (\emph{Notas}, p. 209 a 215). Aristoteles, portanto, admittia o que a sciencia moderna admitte, que não só o \emph{som} é um \emph{movimento}, mas que o são tambem o \emph{calor}, os \emph{cheiros}, e os outros sensiveis. E, comtudo, o philosopho grego admittia, ao mesmo tempo, que o \emph{som}, a \emph{luz} e os outros sensiveis \emph{proprios} são verdadeiramente \emph{objectivos}, e esta \emph{objectividade} era, para elle, a base de todo o conhecimento intellectual. Haverá nisto uma contradicção? Vamos demonstrar que não ha.\par O \emph{movimento}, de que se trata na questão, não é um movimento \emph{abstracto}, mas é um movimento \emph{concreto}, porque é o movimento de uma substancia \emph{extensa} e \emph{activa}. Ora o movimento de uma substancia \emph{extensa} --- não é meramente \emph{local} ou \emph{accidental}, capaz de \emph{transportar} de um logar para outro um phenomeno de luz, de som, etc., --- mas é tambem \emph{essencial} ou \emph{constitutivo}, capaz de \emph{produzir} o proprio phenomeno. Por quanto, devemos distinguir varias especies de movimentos ou de mudanças. Diz Aristoteles (\emph{De nat. auscult.}, III, c. 1): « O movimento não existe fóra das coisas; porque tudo o que se muda, muda-se --- ou emquanto á \emph{substancia} (geração), --- ou emquanto á \emph{quantidade} (augmento), --- ou emquanto á \emph{qualidade} (alteração), --- ou emquanto ao \emph{logar} (translação) ». A distincção entre o movimento puramente \emph{local} e o movimento \emph{essencial} é capitalissima. O \emph{local} é \emph{extenso} e só produz \emph{deslocação}, e consiste na passagem continua da potencia para o acto, como se vê nas oscillações de um pendulo. O \emph{essencial} é mais \emph{intimo}; as vibrações moleculares, dotadas de uma finura e rapidez prodigiosa, alteram os corpos no seu modo de ser, nas suas propriedades physicas e chimicas, e produzem nelles os phenomenos da luz, do calor, da electricidade, etc. Portanto as qualidades sensiveis são um producto do \emph{movimento}, mas do movimento \emph{essencial}, e comtudo e por isso são \emph{objectivas}. --- Como tambem o movimento de uma substancia \emph{activa} não se pode conceber, sem que se conceba, ao mesmo tempo, o effeito d'essa actividade. Uma substancia \emph{activa}, que se move, opera e nada produz, é uma contradicção. E', pois, essencial a uma energia ou actividade a producção de um phenomeno, e esta producção especifica a propria energia. Por isso --- uma energia, desenvolvendo-se, produzirá o phenomeno da extensão continua, --- uma outra produzirá o phenomeno da côr ou do som; e estas differentes energias serão especificadas pelas formas, que tiverem revestido. Se não ha movimento sem um producto, é claro que um movimento, separado, por abstracção, do seu producto, é sempre um \emph{acto incompleto}. Sendo o movimento aperfeiçoado e especificado pelo phenomeno produzido, Aristoteles com muita razão conclue que haverá tantas especies de movimentos, quantas forem as especies dos\pgn{256} seres ou das mudanças nelles realisadas: « Quocirca, motus et mutationis species tot sunt, quot entis » (\emph{De nat. aus.}, III, c. 1). Logo os sensiveis \emph{proprios} differem especificamente entre si pela \emph{sua essencia}, independentemente da percepção dos nossos sentidos. A \emph{luz} é o resultado \emph{especifico} da acção \emph{luminosa}; o \emph{som} é o resultado \emph{especifico} da acção \emph{sonora}; o \emph{extenso continuo} é o resultado \emph{especifico} da \emph{força de extensão}. Os sentidos limitam-se a perceber e constatar este resultado. Diz Aristoteles (\emph{Met.}, IV, 14): « As \emph{qualidades sensiveis} são as \emph{modificações} das substancias postas em movimento, como o \emph{calor}, o \emph{frio}, a \emph{alvura}, etc. A palavra \emph{qualidade} indica a \emph{modificação} das coisas que se movem, emquanto se movem, e as \emph{differenças} dos movimentos ».\par Esta é a conclusão dos grandes mestres. Será ella contraria ás sciencias experimentaes? Não; ellas confessam ignorar o modo por que os diversos movimentos são especificados. --- Será contraria specialmente aos principios da \emph{physica} e da \emph{chimica}? Tambem não; a these da \emph{objectividade} dos sensiveis não altera nenhum theorema, não destroe nenhuma formula ordinaria.\par Os adversarios objectam muitos factos como inconciliaveis com a \emph{objectividade} dos \emph{sons} e das \emph{córes}. Estes factos, dizem elles, provam que todos os nossos sentidos, com excepção do \emph{tacto} (e por isso a \emph{extensão}, objecto do \emph{tacto}, é o unico sensivel \emph{objectivo}), se enganam emquanto percebem não o que é real e objectivo, mas só as modificações do sujeito; e por isso, concluem, o officio dos nossos sentidos reduz-se á percepção das sensações \emph{subjectivas}, que depois transportamos fóra de nós e transformamos em \emph{objectivas} por uma certa illusão natural.\par Não podemos analysar detidamente os factos apontados como inconciliaveis com a \emph{objectividade} do \emph{som} e da \emph{côr}. Não falta quem o tenha feito, e com excellente resultado. Remettemos o leitor ás duas magnificas obras, acima recommendadas. --- Não podemos, porém, deixar de fazer umas observações, que servem para uma refutação \emph{generica} da difficuldade. Antes de tudo, repetimos o que foi provado na \emph{Logica}, isto é, que nenhum sentido se engana, quando se refere ao seu objecto \emph{proprio} e se acha em condições \emph{normaes}. Pode dar-se e dá-se effectivamente um engano, quando se trata de um objecto \emph{improprio}, ou quando a condição do sentido é \emph{anormal}; mas, nesse caso, o erro é apenas \emph{accidental}, e é absurdo concluir que o sentido se engana com relação ao objecto \emph{proprio} e no estado \emph{normal}, só porque se engana com relação ao objecto \emph{improprio} e em condições \emph{anormaes}. Diz S. Thomaz: « Nulla potentia cognoscitiva deficit a cognitione \emph{sui} objecti, nisi propter aliquem defectum aut corruptionem » (\emph{C. Gentes}, III, c. 107). --- Os adversarios confundem, geralmente, a percepção \emph{externa} dos sentidos, que attinge a realidade \emph{externa}, com a percepção \emph{interna} do \emph{sentido commum}, que percebe a realidade \emph{interna}, isto é, as impressões dos nossos orgãos; e estas duas percepções, interna e externa, confundem-nas com a \emph{imaginação}, que fabrica um mundo phantastico. E assim, e só assim, poderam reduzir estas tres especies de phenomenos, tão differentes entre si no mechanismo e no alcance, á mesma formula, chamando-os: \emph{emoções ou modificações do eu}, \emph{percebidas e interpretadas pela intelligencia}! --- Quando dizem que \emph{todos} os sentidos estão sujeitos a illusões, á excepção do \emph{tacto}, dizem uma falsidade; porque tambem o \emph{tacto} tem as suas illusões, como os outros sentidos, e portanto, se\pgn{257} os outros sentidos não merecem confiança, não a merece o \emph{tacto}. --- Além d'isso, não é exacto o que dizem da \emph{extensão}, que seria o unico sensivel \emph{real}, porque objecto do \emph{tacto}. A \emph{extensão} é um objecto \emph{commum}, porque se refere á \emph{vista} e ao \emph{tacto}. Ora o objecto \emph{commum} (como tambem o \emph{concomitante}) não pode ser percebido, senão emquanto são percebidos antecedentemente os sensiveis \emph{proprios}: « Sensibilia \emph{per accidens} non apprehenduntur, nisi prius apprehendantur sensibilia \emph{propria}; nec sensibilia \emph{communia} apprehenduntur, nisi apprehendantur sensibilia \emph{propria}; numquam \emph{visus} apprehendit \emph{magnitudinem} aut \emph{figuram}, nisi in quantum apprehendit \emph{coloratum} » (S. Th., \emph{De Anima}, II, 13). Portanto, a \emph{extensão} não pode ser percebida, pelo menos \emph{adequadamente}, se não for percebida primeiro a côr d'essa extensão. Ora, como a \emph{côr} é percebida pela \emph{vista}, se esta não fornece senão elementos \emph{subjectivos}, segue-se que o conceito \emph{adequado} da \emph{extensão} será composto tambem de elementos \emph{subjectivos}, e assim a \emph{objectividade} da \emph{extensão}, apregoada pelos adversarios, ficará destruida. Nem se diga que o \emph{tacto}, só por si, pode perceber a extensão; porque o \emph{tacto}, se não fôr auxiliado pela \emph{vista}, não é capaz de nos dar uma idea \emph{completa} da \emph{extensão}: o que os proprios adversarios reconhecem. Tal é o resultado de concessões imprudentes. Não ha meio termo nesta questão: ou \emph{objectivistas} em tudo, ou \emph{subjectivistas} em tudo!}.

%%>>> t158

\pg{250} \pg{251} \tese{158}{A percepção dos objectos externos é directa ou immediata.}

a) As faculdades perceptivo-sensitivas \emph{externas}, determinadas ao acto pela imagem \emph{impressa}, percebem os objectos \pg{252}  \emph{externos}, isto é, os sensiveis \emph{proprios}, que são a côr, o som, a figura, e outras qualidades singulares e concretas, entre as quaes se acha como que envolvida a substancia. Ora esses \pg{253}  sensiveis \emph{proprios} são percebidos pelas respectivas faculdades de um modo \emph{directo} ou \emph{immediato}. Por quanto, a acção do \pg{254}  objecto e a paixão do sujeito não são dois actos distinctos, mas constituem um e o mesmo acto, que tem o objecto por \pg{255}  principio e o sujeito por termo. Constituindo a acção do objecto e a paixão do sujeito um o o mesmo acto, o sujeito, para perceber \pg{256}  o objecto, não precisa de intermediario, mas, pelo facto de ter o mesmo objecto em si mesmo, percebe-o \emph{directa} ou \pg{257}  \emph{immediatamente}. Logo a percepção dos objectos \emph{externos} é \emph{directa} ou \emph{immediata} \footnote{Desenvolvamos o argumento exposto no texto. --- O sentido está para o sensivel, como o paciente para o agente: « \emph{Sensus comparatur ad sensibile}, \emph{sicut patiens ad agens} » (S. Th., \emph{De Verit}., q. 26, a. 3). Ora o agente, actuando no paciente, une-se a este pela acto; de modo que o agente e o paciente, embora distinctos na \emph{substancia}, se acham unidos no \emph{mesmo acto}. Um e o mesmo é, portanto, o acto, que o agente produz e o paciente recebe, embora no agente se encontre como no seu \emph{principio} e no paciente como no seu \emph{sujeito}. « \emph{Operatio} rei sensibilis atque sensus \emph{una} quidem et \emph{eadem est}; et eorum \emph{ratio non est eadem} » (Arist., \emph{De Anima}, III, c. 2). Por outras palavras: a acção e a paixão são um e o mesmo acto, porque é o acto do agente que informa a potencia do paciente. Por isso diz S. Thomaz que o sentido \emph{em acto} (emquanto é \emph{actuado} pelo sensivel) é o proprio sensivel \emph{em acto} (emquanto é \emph{actualmente} percebido), porque do sentido e do sensivel resulta um unico ser (na ordem \emph{intencional}), composto de potencia e de acto (o \emph{sentido} é a \emph{potencia}, o \emph{sensivel} é o \emph{acto}): « \emph{Sensus in actu est sensibile in actu}, \emph{quia ex utroque fit unum}, \emph{sicut ex potentia ex actu} » (\emph{Sum. Th.}, p. I, q. 55, a. 1 ad 2). --- Sendo o mesmo acto \emph{commum} ao agente e ao paciente, é claro que este não precisa de sahir de si mesmo para o perceber. Portanto o paciente percebe \emph{dentro de si mesmo} o acto do agente externo, e percebe-o como elle é, isto é, como um acto recebido e não produzido, como um acto que foi acolhido no \emph{eu} mas que partiu do \emph{não-eu}; e, por uma reacção natural e espontanea, vira-se para a origem do acto e attribue ou restitue ao agente o que lhe pertence e que é o proprio acto. --- O mesmo dar-se-hia com a cera,\pgn{258}se esta tivesse consciencia. Recebida a impressão de um sinete, a cera perceberia necessariamente a figura gravada nesse sinete; porque o acto do agente é identico ao do paciente. --- Mas o acto é a expressão, a manifestação, embora \emph{inadequada}, do agente, é alguma coisa d'elle; porque o agente não pode operar e não opera senão em conformidade com a natureza ou com a sua forma. Por isso, se o acto é luminoso, figurado, extenso, o paciente attribuirá essas qualidades sensiveis ao agente externo, e perceberá, implicitamente, a propria \emph{substancia} do agente, impenetravel aos sentidos, mas envolvida nessas qualidas. --- Logo, se acção e a paixão constituem um unico acto, --- se a acção do agente se encontra no paciente, --- se o paciente percebe o agente pela acção, que este exerce naquelle, --- deve concluir-se que o paciente distingue e percebe \emph{directa} ou \emph{immediatamente}, sem intermediario, a acção do agente; e por isso a percepção, que o paciente ou a faculdade sensitiva adquire do agente ou do corpo externo, é \emph{directa} ou \emph{immediata}. --- E' verdade que os agentes externos não actuam sempre \emph{immediatamente} nos sentidos. Ás vezes é a ultima \emph{camada} do meio interposto, á qual os corpos communicam as suas vibrações luminosas ou sonoras, que actua \emph{immediatamente} na retina ou no nervo acustico; o que já ensinàra S. Thomaz nas seguintes palavras: « \emph{In aliis sensibus immutatio medii est causa quod immutetur sensus, non autem in tactu} » (\emph{De Anima}, II, l. 23). Mas tambem n'estes casos a percepção deve dizer-se \emph{directa} ou \emph{immediata}, porque entre o objecto e o sentido não ha outro intermediario a não ser o \emph{acto}, que parte de um e é recebido no outro. --- Como se vê, o argumento exposto prova tambem a these defendida contra os \emph{subjectivistas}.\par Dizem: No argumento adduzido suppõe-se como provado e certissimo que a acção é \emph{semelhante} ao agente. Ora não é bem assim. O principio --- \emph{omne agens agit sibi simile} --- não quer dizer que o paciente se torna semelhante ao agente, mas indica apenas que o agente opera de um modo \emph{proporcionado} á propria natureza.\par Resposta. O sentido obvio d'esse principio é que o agente, em virtude do seu acto, tende, por quanto depende d'elle, a imprimir no paciente a propria forma, não real mas \emph{intencional}. E nada mais logico. Operando o ente pela sua \emph{forma}, o seu acto não pode deixar de ser a expressão, a manifestação da mesma forma; assim um corpo \emph{quente}, isto é, dotado actualmente da \emph{forma de calor}, tende, com as suas irradiações calorificas, a aquecer os corpos circumstantes, e assim o acto exprimirá e manifestará a mesma forma. Repetimos con S. Thomaz: « \emph{De natura agentis est, ut agens sibi simile agat, quum unumquodque agat secundum quod actu est} » (\emph{C. Gent.}, I, 29). --- Sendo o acto a expressão, a manifestação do agente, segue-se que receber o \emph{acto} é receber a \emph{semelhança} do agente, ou tornar-se \emph{semelhante} ao agente. --- Os adversarios excluem a \emph{semelhança} e admittem uma \emph{proporção}, como se estas duas coisas fossem incompativeis. Mas essas duas coisas, longe de serem incompativeis, relacionam-se por tal maneira, que uma suppõe a outra. Por isso mesmo que o agente opera de um modo \emph{proporcionado} á sua natureza, deve imprimir a propria \emph{semelhança} no sujeito, que recebe a\pgn{259}acção. Porquanto não se pode operar d'esse modo, sem que a natureza \emph{se reflicta} no acto e por isso no effeito, isto é, sem que o effeito seja \emph{semelhante} ao agente; assim, por isso mesmo que o sol opera de um modo \emph{proporcionado} á sua natureza, reproduz a propria semelhança nos corpos sujeitos á sua acção, emquanto estes ficam acquecidos e illuminados. --- Não se pode nem se deve eliminar do tratado da percepção \emph{sensitiva} a theoria da \emph{assimilação}, como pretendem alguns \emph{neo-escholasticos}, só porque ainda não descobriram uma solução satisfactoria para certas difficuldades.}.

b) As faculdades perceptivo-sensitivas \emph{externas} devem sempre perceber alguma coisa de um modo \emph{directo} ou \emph{immediato}. Essa coisa --- ou é a imagem \emph{impressa}, --- ou a imagem \emph{expressa}, --- ou a modificação \emph{subjectiva} do orgão, --- ou o \pg{258}  objecto \emph{externo}; porque não se pode conceber outra hypothese. Ora, conforme a experiencia \emph{interna} e o \emph{senso commum} nos attestam, o que essas faculdades \emph{directa} ou \emph{immediatamente} percebem --- não é a especie \emph{impressa}, --- nem a especie \emph{expressa}, --- \pg{259}  nem a modificação \emph{subjectiva} do orgão, --- mas é o proprio objecto \emph{externo}. Logo a percepção dos objectos \emph{externos} é \emph{directa} ou \emph{immediata} \footnote{Dizemos no texto que, segundo o attestado da \emph{experiencia interna} e do senso \emph{commum}, é o objecto \emph{externo} que se percebe \emph{directa} ou \emph{immediatamente}, e não a \emph{imagem impressa}, nem a \emph{imagem expressa}, nem a \emph{impressão organica}. Desenvolvamos tambem este ponto.\par \emph{a}) Não é a \emph{imagem impressa}. Na verdade, a imagem \emph{impressa} é o que a acção do agente produz no paciente, e por isso se chama tambem \emph{paixão}. Ora o que os sentidos percebem de um modo \emph{directo} ou \emph{immediato} não é a \emph{paixão}, mas é a \emph{acção}. Assim, quando um objecto de figura \emph{concava} actua nos sentidos e produz nelles uma \emph{paixão} de figura \emph{convexa}, o que a faculdade percebe immediatamente não é a figura \emph{convexa}, mas é a \emph{concava}, isto é, não é a \emph{paixão} mas a \emph{acção}. Como tambem, quando a acção luminosa de um objecto imprime na retina a imagem \emph{invertida} desse objecto, os olhos não vêm a imagem \emph{invertida}, mas vêm o objecto na sua posição \emph{normal}, isto é, não vêm a \emph{paixão} mas a \emph{acção}. Portanto a \emph{acção} é o objecto que se vê \emph{directa} ou \emph{immediatamente} (est id \emph{quod} percipitur); emquanto a \emph{paixão}, ou a imagem \emph{impressa}, é o \emph{meio}, \emph{pelo qual} o \emph{objecto} se percebe (est id \emph{quo} objectum percipitur), é apenas uma \emph{condição}, embora indispensavel, para a união \emph{immediata} do agente com o paciente, e para a percepção da acção do agente no paciente. --- E' uma lei psychologica, que não possamos conhecer o \emph{eu} senão pela sua opposição com o não-eu, e viceversa. Mas, embora a acção só se perceba na \emph{paixão}, comtudo é sempre a \emph{acção} que se percebe \emph{immediatamente} na \emph{paixão}, onde aquella se reflecte. --- Esta é a doutrina do Angelico Mestre. Diz elle: « Species (\emph{impressa}), quae est in visu, non est id \emph{quod} videtur, sed \emph{quo} visus videt; \emph{quod} autem videtur est \emph{color}, \emph{qui est in corpore} » (\emph{De Anima}, III, 1. 9). E noutra parte: « Species (\emph{impressa}) \emph{abstrahitur} (h. e. \emph{derivatur}) a rebus ipsis, et \emph{per} eam cognitio sensus \emph{continuatur} ad \emph{ipsas res sensibiles} » (\emph{De Verit.}, q. II, a. 6). E ainda: « Species (\emph{impressa}) non est illud \emph{quod} sentimus, sed magis illud \emph{quo} sensus sentit » (\emph{Sum. Th.}, p. I, q. 85, a. 2). --- Portanto a imagem \emph{impressa} pelo agente no paciente --- não é o \emph{termo}, mas o \emph{principio} da percepção, --- não é um \emph{signal objectivo}, que deva ser \emph{anteriormente} percebido, mas é o \emph{meio} para unir o agente com o paciente.\par \emph{b}) Não é a especie \emph{expressa}. A especie \emph{expressa} (a qual corresponde \emph{analogicamente} ao \emph{verbo} da percepção \emph{intellectual}) é o \emph{proprio acto perceptivo}. Por isso, se um sentido percebesse a sua especie \emph{expressa}, perceberia o seu acto. Ora os sentidos não podem perceber o seu acto, porque são destituidos de reflexão \emph{perfeita}, e muito menos podem percebel-o de um modo \emph{directo} ou \emph{immediato}; porque o termo \emph{directo} ou \emph{immediato} do acto não é proprio\pgn{260}acto, mas é o \emph{objecto}, para o qual o acto tende por um movimento \emph{essencial} e \emph{primario}, e que \emph{especifica} o mesmo acto. --- S. Thomaz compara a especie \emph{expressa} a um espelho, cujas dimensões não excedem as do objecto, que nelle se reflecte, de modo que as dimensões do espelho e do objecto são perfeitamente eguaes. Diz assim: « Species expressa est tamquam speculum, in quo res cernitur, sed non excedens id, quod in eo cernitur » (\emph{De natura verbi intellectus}). Ora quem fita esse espelho, não vê \emph{primeiramente} o \emph{espelho} e \emph{em seguida} o \emph{objecto}, mas vê \emph{immediata} e \emph{primeiramente} o \emph{objecto}, porque o \emph{espelho} occulta-se na sua propria reflexibilidade, e limita-se, por sua natureza, a mostrar o proprio objecto.\par c) Não é a \emph{impressão organica}. Esta impressão pode considerar-se --- ou emquanto é mera affecção \emph{subjectiva} do paciente, --- ou emquanto é uma paixão \emph{representativa} da acção do agente. --- Se se considera emquanto é uma affecção \emph{subjectiva} do paciente, não é o que os sentidos percebem de um modo \emph{directo} ou \emph{immediato}; porque a experiencia \emph{interna} nos attesta que, quando a ponta de uma agulha nos fere um dedo, \emph{primeira} ou \emph{immediatamente} percebemos a ponta da agulha, e só em seguida sentimos a dôr. --- Se se considera emquanto é uma paixão \emph{representativa} da acção do agente, então identifica-se com a imagem \emph{impressa}, e já vimos que não é \emph{esta} o objecto \emph{directo} ou \emph{immediato} da percepção sensitiva.\par Objectam: A percepção sensitiva não se pode chamar \emph{directa} ou \emph{immediata}, porque entre a faculdade e o objecto \emph{medeia} a imagem \emph{impressa}.\par Resposta. A existencia da imagem \emph{impressa} entre a faculdade e o objecto tornaria a percepção \emph{indirecta} ou \emph{mediata}, se essa imagem fosse o \emph{termo} da percepção, isto é, se fosse um \emph{meio objectivo}, que devesse ser percebido primeiramente em si mesmo e depois levasse a faculdade a perceber o objecto externo. Mas não é assim. A imagem \emph{impressa}, como vimos, não é o \emph{termo} da percepção, não é um meio \emph{objectivo}, mas é o meio \emph{formal}, que leva \emph{immediatamente} a faculdade a perceber o objecto externo, sem que elle seja anteriormente percebido, e por isso não torna \emph{indirecta} ou \emph{mediata} a percepção. --- Como adverte S. Thomaz, dá-se um \emph{triplice} meio de percepção: o meio \emph{sub quo}, o meio \emph{quo}, o meio \emph{in quo}. O meio \emph{sub quo} é o que, pela sua acção, torna cognoscivel \emph{em acto} o objecto; tal é a luz do sol com respeito ás côres. O meio \emph{quo} é o que, pela sua natureza, une actualmente o objecto com a faculdade e especifica o acto da propria faculdade, e por isso não é o \emph{termo}, mas o principio \emph{formal} da percepção; tal é a especie ou imagem \emph{impressa}, sensivel ou intelligivel. O meio \emph{in quo} é o que, depois de ter sido conhecido em si mesmo, leva a faculdade a conhecer um outro objecto; tal é o \emph{effeito}, que antes se conhece em si mesmo e em que depois se percebe a \emph{causa}, e tal é o espelho, cuja visão leva á visão dos objectos, que nelle se reflectem. Ora é evidente --- que nem o meio \emph{sub quo}, nem o meio \emph{quo}, tornam \emph{indirecta} ou \emph{mediata} a percepção, mas só o meio \emph{in quo}, --- e que a imagem \emph{impressa} é o meio \emph{quo}. Ouçamos o S. Doutor: « \emph{Immediate} enim dicitur aliquis videre lapidem, quamvis eum \emph{per speciem} ejus in oculo receptam et \emph{per lumen} videat; quia visus non \emph{fertur} in haec media, tamquam \emph{in visibilia}; sed \emph{per haec media} fertur in \emph{unum visibile}, \emph{quod est extra oculum},\pgn{261}Sed \emph{tertium} medium facit visionem \emph{mediatam}; visus enim primo fertur in speculum \emph{sicut in visibile}, quo mediante accipit speciem rei visae in speculo » (\emph{Quodl.} 8, a. 1; \emph{De Anima}, a. 7 ad 8).\par Sendo \emph{directa} ou \emph{immediata} a percepção dos corpos externos, segue-se que a existencia d'elles não precisa de ser demonstrada pelo raciocinio; o que se vê, não se demonstra. A demonstração é necessaria, quando um objecto não se vê em si mesmo, e por isso, para o conhecermos, devemos partir de um objecto anteriormente conhecido; porque a demonstração consiste precisamente na transição de um objecto conhecido para um outro desconhecido. Ora já vimos que a imagem \emph{impressa} não é um meio \emph{objectivo}, que deva ser anteriormente percebido, mas é um meio \emph{formal}, que leva \emph{directa} ou \emph{immediatamente} a faculdade a perceber o objecto, que ella --- a imagem --- representa. Diz S. Thomaz: « \emph{Cognoscere aliquid per alterum contingit dupliciter}. Uno modo sicut cognoscere \emph{unum cognitum} per \emph{aliud cognitum}, ita quod sit \emph{distincta} cognitio \emph{utriusque}; sicut homo per principia cognoscit conclusionem, \emph{seorsum considerando utrumque}. Alio modo sicut cognoscitur aliquid cognitum \emph{per speciem}, \emph{qua} cognoscitur; ut videmus lapidem per speciem lapidis, quae est in oculo. Primo igitur modo cognoscere unum per alterum facit \emph{discursum}, \emph{non autem secundo modo} » (\emph{De Anima}, a. 7, ad 8).\par *\par Aqui devemos fazer umas observações muito importantes, relativas á percepção do \emph{triplice} objecto dos sentidos.\par \emph{a}) O objecto \emph{formal} da percepção sensitiva é constituido, como dissemos, por certas qualidades \emph{activas} dos corpos, chamadas \emph{sensiveis} e percebidas de um modo \emph{directo} e \emph{immediato}. Mas essas qualidades não são entes completos, subsistentes em si; não são \emph{entes}, são entes do ente, isto é, da \emph{substancia}. Por isso, quando dizemos que o objecto \emph{proprio} e \emph{formal}, que os sentidos attingem de um modo \emph{directo} ou \emph{immediato}, é constituido pelas qualidades sensiveis, não entendemos que estas qualidades são o objecto \emph{adequado} da percepção; mas entendemos que, pertencendo a uma substancia corporea a natureza e as qualidades, a imagem \emph{impressa} representa essa \emph{substancia} emquanto ás \emph{qualidades}, e não emquanto aos \emph{principios constitutivos} da \emph{natureza}. --- Esta foi sempre a doutrina dos escholasticos. Disseram sempre que os sentidos não percebem as qualidades sensiveis de um modo \emph{abstracto}, como separadas da substancia e existentes em si mesmas, mas as percebem como ellas se encontram no objecto, isto é, sustentadas pela \emph{substancia}. Aristoteles diz sempre: « Non immutat visum \emph{albedo}, sed \emph{album} »; e S. Thomaz repete a cada passo: « Sensus est apprehendere \emph{hoc coloratum} ». Logo o objecto \emph{adequado}, que os sentidos percebem de um modo \emph{directo} ou \emph{immediato}, é a \emph{substancia} como é em si mesma, isto é, o \emph{composto} de \emph{substancia} e de \emph{accidentes}. E nada mais natural. As qualidades sensiveis, quando actuam nos sentidos, actuam como coisas \emph{subsistentes}, e como taes são percebidas. Mas ellas são subsistentes --- não pela subsistencia \emph{propria} --- mas pela subsistencia da \emph{substancia}, a que adherem. Por isso a percepção das qualidades sensiveis importa necessariamente a percepção de um attributo, que é proprio da \emph{substancia} e que é a \emph{subsistencia}, embora a \emph{natureza} da substancia escape\pgn{262}ao alcance dos sentidos. --- Accrescente-se que as qualidades sensiveis não impressionam os sentidos como agentes \emph{separados}, mas como \emph{existentes} na \emph{substancia}; porque o ente que não tem subsistencia propria não tem acto proprio (\emph{operari sequitur esse}). Por isso quem opera é a propria \emph{substancia}, ainda que em razão das suas qualidades (a substancia, por si, não pode impressionar os sentidos). Assim a luz, o calor, a gravidade, etc., não operam por conta propria; são principios activos, pelos quaes as substancias operam umas nas outras. E se a acção do agente deriva da \emph{substancia}, embora esta opere em razão das qualidades, a \emph{reacção} do paciente deve ter por termo --- não os qualidades \emph{sensiveis}, como separadas e existentes em si; --- mas a \emph{substancia}, emquanto é revestida d'essas qualidades. Assim não é a \emph{dureza}, a \emph{alvura}, etc., que os sentidos percebem, mas o \emph{corpo duro}, \emph{alvo}, etc. E a imagem ou especie \emph{expressa} representará não a \emph{acção abstracta} do agente, mas o \emph{composto} do agente e da acção. --- Todavia, as qualidades \emph{sensiveis} são percebidas \emph{formal} e \emph{explicitamente}, emquanto são reproduzidas na imagem \emph{expressa} como são em si mesmas; ao passo que a \emph{substancia} é percebida só \emph{implicitamente}, emquanto os seus \emph{constitutivos}, genericos ou especificos, não são reproduzidos naquella imagem. --- Do facto de ser apenas \emph{implicita} a percepção da \emph{substancia}, segue-se que se, por milagre, as qualidades sensiveis fossem separadas do proprio sujeito, a percepção deixaria de incluir a \emph{substancia}, sem, comtudo, soffrer uma intrinseca mudança. Temos d'esta verdade um unico, mas certissimo exemplo no adoravel Sacramento da Eucharistia. A \emph{percepção} das especies sensiveis do pão e do vinho é a \emph{mesma} antes e depois da consagração, e comtudo, depois da consagração, as especies não estão adherentes á substancia do pão e do vinho, como estavam antes. --- D'isto se vê o erro dos que, --- com \emph{Locke}, \emph{Berkeley}, \emph{Reid}, \emph{Mallebranche}, \emph{Rosmini}, \emph{Balmes}, \emph{Kant} e com todos os adversarios da \emph{objectividade} da \emph{substancia}, --- affirmam que a percepção sensitiva externa é \emph{abstracta}, e não \emph{concreta}.\par \emph{b}) O objecto \emph{per accidens}, como é a \emph{substancia}, pelo facto de estar intimamente unido com as qualidades \emph{sensiveis}, é percebido com o \emph{mesmo acto} e do mesmo modo \emph{directo} ou \emph{immediato}. Isto só se verifica quando o objecto \emph{per accidens} é a \emph{substancia}. Porquanto, se esse objecto é uma outra qualidade sensivel, não se segue que seja percebida pelo mesmo sentido, que percebe a primeira. Assim, sendo a \emph{doçura} de uma maçã um \emph{sensivel per accidens} com relação á \emph{vista}, é claro que esta, se percebe a \emph{côr vermelha} da maçã, não percebe a \emph{doçura}. --- Pode chamar-se tambem objecto \emph{per accidens} dos sentidos uma propriedade, que está unida com o sensivel \emph{proprio} e com este é percebida --- não pelo sentido (porque não está representada na imagem \emph{impressa}) mas por uma faculdade \emph{superior}; assim, vendo um individuo humano, lembro-me de um acto bom, praticado por elle. Neste caso, o nexo entre o sensivel \emph{per se} e o \emph{accessorio} não se funda na união \emph{physica} dos dois objectos, mas consiste na apprehensão do facto precedente e do objecto presente, e a comparação é feita pela \emph{cogitativa}, á qual pertence relacionar as noções \emph{individuaes}, como a intelligencia relaciona as noções \emph{universaes}. (Cf. S. Th., \emph{De Anima}, II, 1. 18).\par c) Os objectos \emph{communs}, como dissemos, são os que referem a mais de um sentido. Estes objectos são intermedios entre os \emph{proprios} e os \emph{per accidens}.\pgn{263}--- Se o \emph{movimento} e a \emph{figura} não são qualidades \emph{activas}, que modifiquem \emph{directamente} o orgão, \emph{modificam}, comtudo, a impressão causada pelos objectos \emph{proprios}. --- Os objectos \emph{communs} não são representados por uma especie \emph{propria}, exclusivamente propria d'elles; mas são representados na especie ou imagem dos sensiveis \emph{proprios}, emquanto essa imagem é diversamente modificada, isto é, emquanto representa os sensiveis \emph{proprios}, diversamente \emph{modificados}. --- A \emph{extensão} é percebida, emquanto a acção da qualidade sensivel, por ser a de um agente \emph{material}, se exerce de um modo \emph{extenso}. \emph{A figura} percebe-se, emquanto acrescenta á extensão certos e determinados limites. \emph{O numero} é percebido, emquanto as qualidades da mesma ordem, concretisadas em varios individuos semelhantes, impressionam conjunctamente os sentidos. O \emph{movimento} não é percebido em si mesmo, mas nas proprias acções das coisas, que se mudam e que produzem successivas impressões nas diversas partes da faculdade perceptiva. --- Sendo percebidos com os sensiveis \emph{proprios}, os sensiveis \emph{communs} são tambem percebidos de um modo \emph{directo} ou \emph{immediato}.\par *\par Do que deixamos dicto vê-se quão vão é o esforço dos que procuram uma \emph{ponte} entre o objecto externo e o sujeito sensiente. A ponte, como dissemos, serve para unir distancias; mas o objecto e o sujeito não são distantes, porque se acham unidos no mesmo acto, e a especie \emph{impressa} é o laço que os une. Quando eu, tocando numa pedra, sinto uma coisa que me resiste, que abysmo poderá existir entre a minha mão e a pedra? Nenhum. Feita a impressão, a pedra encontra-se na minha mão pela sua resistencia. --- Esta é a razão da convicção universal e constante de que, percebendo pelos sentidos o mundo exterior, não percebemos apenas as nossas modificações subjectivas, mas attingimos \emph{directa} e \emph{immediatamente} alguma coisa propria dos objectos externos, isto é, as suas acções naturaes e por estas a substancia. E' a \emph{realidade}, é o \emph{absoluto} que percebemos, embora essa percepção seja \emph{fragmentada} e exija um trabalho, espontaneo ou reflexo, de reconstrucção do objecto total. --- Só os \emph{subjectivistas} contradizem este consenso unanime do genero humano, e então, se não chegaram á consequencia extrema de negar a propria existencia do mundo corporeo (como fizeram \emph{Berkeley}, \emph{Hume}, \emph{Kant}, etc.), sentiram a necessidade de procurar uma \emph{ponte}, pela qual podessem passar do seu \emph{subjectivismo} para a externa e \emph{objectiva} realidade das coisas. Mas a \emph{ponte} excogitada não serviu para o fim. Não vale a pena referir todas as opiniões dos \emph{subjectivistas} moderados a este respeito. --- Limitamo-nos a referir as de \emph{Descartes}, de \emph{Mallebranche} e de \emph{Reid}.\par *\par \emph{a}) \emph{Descartes}, seguido por \emph{Romagnosi}, disse que a certeza ácerca da real existencia dos corpos se deduz pelo raciocinio, como da causa se deduz o effeito, emquanto deve haver uma causa das nossas sensações, e esta causa só pode ser a real existencia dos corpos. Esta opinião, á qual deram o nome de \emph{inferencia} ou \emph{illação}, foi sustentada nos ultimos tempos por \emph{Cousin}.\par A conclusão de \emph{Descartes} não é logica. Por quanto, se, na theoria dos \emph{subjectivistas}, as sensações não representam de modo algum as realidades objectivas e externas, não é possivel deduzir da existencia das sensações a existencia dos objectos externos; porque, se o effeito nos leva ao conhecimento da causa, a razão é que no effeito se reflecte de algum modo a propria causa. --- Alem disso, se a percepção dos objectos externos devesse attribuir-se a um raciocinio, os animaes, que são destituidos de razão, não poderiam percebel-os. E todavia os animaes percebem esses objectos, porque os procuram, e não poderiam procural-os, se os não percebessem. --- \emph{Descartes} recorre á \emph{veracidade divina} e diz que, embora nós percebamos apenas as nossas sensações e não os objectos externos, comtudo estamos certos de que esses objectos existem, porque foi Deus quem deu a \emph{todos} os homens uma inclinação invencivel para acreditarem na \emph{evidencia} dos nossos sentidos, e Deus não pode enganar. Mas este recurso, além de não ser serio, envolve contradicção. Como sabe o escriptor francez que a inclinação, que nos leva a acreditar na \emph{evidencia} dos sentidos, é \emph{universal} e \emph{invencivel}, e por isso imputavel ao Auctor da natureza? Só o pode saber pela \emph{percepção sensitiva}, que nos põe em communicação com os outros homens, e por isso pelo \emph{veraz testemunho} dos sentidos, que elle não quer reconhecer. --- Alem d'isso, affirma \emph{Descartes} que temos a inclinação para acreditar na \emph{evidencia} dos sentidos. Mas se o testemunho dos sentidos é \emph{evidente}, não precisa de ser demonstrado pelo raciocinio. --- Mais: affirma \emph{Descartes} que Deus não pode permittir que \emph{todos} os sentidos se enganem. Mas se Deus permitte, na opinião de \emph{Descartes}, que se enganam os olhos com relação á côr, porque não poderá permittir que se enganem os outros sentidos com relação a outros objectos? A \emph{veracidade} dos sentidos não é um \emph{direito}, que se reclama, mas é um \emph{facto}, que se deve constatar.\par \emph{b}) \emph{Mallebranche} diz que não podemos estar certos de que os corpos existem, senão pela revelação divina, pela fé em Deus, que nos ensina na S. Escriptura ser Elle o creador do ceu e da terra (\emph{Entretiens}, VI).\par Parece incrivel que um homem de talento, como \emph{Mallebranche}, recorresse a este expediente! Nega-se a \emph{percepção} dos corpos, e depois, em philosophia, \emph{acredita-se} na \emph{revelação} dos corpos! Nega-se o \emph{visivel}, e depois affirma-se o \emph{invisivel}! --- O systema de \emph{Mallebranche} contém o \emph{ontologismo}, e por isso a elle nos havemos de referir no capitulo seguinte.\par \emph{c}) \emph{Reid} encontrou a \emph{ponte} no \emph{cego instincto} da natureza. Disse elle que a nossa intelligencia, quando lhe apparece alguma representação sensivel, é levada, --- por uma \emph{suggestão}, ou \emph{inspiração}, ou \emph{revelação natural}, ou \emph{instincto irresistivel}, --- a affirmar a existencia dos corpos.\par Tambem este recurso do philosopho escocez é vão. Se as sensações são meramente subjectivas, não ha razão por que se deva obedecer a esse \emph{instincto cego} e concluir para a \emph{objectividade} dos corpos externos. --- A convicção do genero humano, relativa á existencia real dos corpos, longe de ser o resultado de uma fé cega, está baseada na \emph{evidencia} da percepção \emph{immediata} dos corpos. --- Todos acreditam na existencia e objectividade do mundo exterior, porque o vêm com os olhos, e o tocam com as mãos. Um acto de fé cega não explica a convicção do bom senso; nega-a e destroe-a.\par *\par Não admira, pois, que, perante tão absurdos desvarios e tão perniciosas consequencias, muitos physiologistas modernos tenham sentido a necessidade de voltar á doutrina dos escholasticos, muito antiga e sempre nova, como\pgn{265}a verdade, e comecem a reconhecer e admittir a percepção \emph{directa} ou \emph{immediata} dos objectos externos. Citaremos, entre muitos, dois: um italiano, outro francez. O italiano é \emph{Bonucci}, o francez é \emph{Bergson}.\par \emph{Bonucci} reconhece e declara lealmente que os olhos não vêem as \emph{imagens da retina}, \emph{mas as coisas externas}. --- E assim, depois de tantas censuras feitas aos escholasticos, por elles ensinarem que o objecto \emph{directo} das nossas faculdades sensitivas é constituido pelas \emph{coisas externas} e \emph{não pelas} suas \emph{imagens}; depois de tantos esforços dos philosophos e dos physiologistas para explicarem --- como é que o cerebro refere as sensações organicas ás varias partes do corpo e não a si mesmo, --- como por meio de imagens \emph{invertidas} se vêem os objectos na sua posição real, --- como por meio de imagens pequenas se percebem objectos grandes, --- como por meio de imagens adherentes á retina ou ao cerebro se vêem objectos distantes, --- como por meio de imagens planas se vêem objectos solidos, --- como por meio da duplice imagem existente nos dois olhos se vê um só objecto, --- depois de tudo isto é verdadeiramente admiravel que \emph{as mais agudas} intelligencias, entre as quaes \emph{Bonucci}, tenham sido obrigadas a confessar \emph{que a faculdade visual não percebe as imagens da retina}, \emph{mas as coisas externas}. (Bonucci, \emph{Della visione delle cose esterne}, 1862, citado por \emph{Panizza} no livro --- \emph{La fisiologia del sistema nervoso}, Roma, 1889, p. 245).\par \emph{Bergson}, professor no Collegio de França diz: « Nós distinguimos na nossa percepção, ao mesmo tempo, um estado da nossa consciencia e uma \emph{realidade independente} de nós... Se, ao contrario, a minha consciencia, longe de distinguir em mim uma acção extranha, só podesse perceber o \emph{eu} e as modificações do eu, não se poderia explicar de modo algum a sensação da exterioridade e da presença do \emph{não-eu}, que experimentamos ». \emph{Matière et memoire}, p. 227, 244).}.

%%>>> t159

\pg{260} \tese{159}{Natureza e propriedades da percepção sensitiva externa.}

a) Natureza da percepção sensitiva externa. --- A percepção sensitiva é, pois, um phenomeno complexo, que resulta de dois

\pg{261}  elementos: da \emph{impressão} e da \emph{percepção} propriamente dicta. A \emph{impressão} é a modificação do sujeito sensiente, causada pela \pg{262}  acção do objecto externo. A \emph{percepção} é um acto vital, pelo qual o sujeito sensiente, reagindo sobre a acção que o impressionou \pg{263} , \emph{exprime} em si a especie ou semelhança da mesma acção e, nesta expressão, percebe o objecto externo. Resultando \pg{264}  da acção do objecto no sujeito e da reacção do sujeito no objecto, o phenomeno da percepção sensitiva é, ao mesmo \pg{265}  tempo (embora sob differentes aspectos), \emph{passivo} e \emph{activo}, \emph{subjectivo} e \emph{objectivo} \footnote{A faculdade sensitiva é \emph{passiva}, emquanto não pode passar da potencia para o acto, senão receber uma impressão do objecto externo. Mas não é um movel meramente passivo. Tem ella uma força de operar; desenvolve uma \emph{actividade}, que lhe é propria, que deriva da sua natureza. Esta actividade, escondida no paciente, recebe um \emph{estimulo}, uma \emph{excitação} pela acção do agente. Recebido o \emph{estimulo}, a faculdade reage, e, \emph{exprimindo} em si a imagem \emph{impressa}, percebe o proprio objecto. Ha, pois, duas phases na percepção sensitiva --- uma \emph{passiva}, outra \emph{activa}. A phase \emph{passiva} é constituida pela \emph{impressão}; a phase \emph{activa}, pela \emph{percepção} propriamente dicta. Estas duas phases são indicadas por S. Thomaz nas seguintes palavras: « Sentire, quantum ad receptionem speciei sensibilis, nominat \emph{passionem}; sed quantum ad actum, consequentem ipsum sensum, perfectum per especiem, nominat \emph{operationem} » (In I \emph{Sent}., Dist. 40, q. 1, a. 1 ad 1). --- A \emph{impressão}, portanto, e a \emph{percepção}, embora estreitamente ligadas entre si, differem uma da outra pelas seguintes razões: --- 1.ª) a \emph{impressão}, por ser o resultado da acção do objecto no orgão, é \emph{passiva}; ao passo que a \emph{percepção}, por ser o effeito da reacção da faculdade no objecto e por isso um acto vital, é \emph{activa}; --- 2.ª) a \emph{impressão} é um phenomeno physico, accessivel á experiencia externa, e pode deixar de ser acompanhada pela \emph{percepção} (quando o orgão se não encontra no estado normal); a \emph{percepção} é um facto vital, psychologico, subjectivo,\pgn{266}accessivel unicamente á consciencia \emph{interna}; --- 3.ª) a \emph{impressão} é multiplice pela multiplicidade das partes, de que consta o objecto; a \emph{percepção} é \emph{simples}, e isto prova a existencia e a simplicidade de um principio sensiente; --- 4.ª) a \emph{impressão} é proporcionada e semelhante á natureza do \emph{agente}, que é o objecto externo; a \emph{percepção} é proporcionada e semelhante á natureza do \emph{paciente}, que é a faculdade perceptiva.\par A \emph{impressão}, como tambem a \emph{percepção}, exige a \emph{união} do objecto com a faculdade; porque, na \emph{impressão}, une-se o objecto com a faculdade; e, na \emph{percepção}, une-se a faculdade com o objecto. --- Para essa união são necessarias varias condições: umas referem-se ao \emph{objecto sensivel}, e outras ao \emph{orgão sensitivo}. --- As condições relativas ao \emph{objecto sensivel} resumem-se nesta: o objecto deve estar \emph{presente} ao orgão, para que possa imprimir neste a sua imagem. Esta presença pode ser --- ou \emph{immediata}, como é a dos objectos do \emph{olfacto}, do \emph{gosto}, do \emph{tacto}, --- ou \emph{mediata}, como é dos objectos dos dois sentidos mais nobres, que são a \emph{vista} e o \emph{ouvido}. --- As condições relativas ao \emph{orgão sensitivo} são as seguintes: --- 1.ª) o orgão deve ser \emph{animado}, \emph{são} e convenientemente \emph{disposto} e \emph{applicado}; o que é evidente; --- 2.ª) o nervo sensitivo deve ter uma sufficiente excitabilidade, produzida pelo \emph{agente nerveo}; porque, sem essa excitabilidade, não pode haver a reacção, necessaria para a percepção do objecto; --- 3.ª) a faculdade sensitiva, para que possa receber em si a especie ou semelhança do objecto sensivel, não deve ter \emph{de modo algum} as qualidades do proprio objecto, --- ou não as deve ter no \emph{mesmo grau}; aliás será impossivel a percepção; --- 4.ª) a faculdade deve ser \emph{proporcionada} ao objecto; por isso, se o sensivel é \emph{commum}, é indispensavel o exercicio de \emph{todos} os sentidos, necessarios para a percepção d'esse sensivel, aliás a percepção não deixará de ser erronea, como p. e. é erronea a percepção dos que só empregam a \emph{vista} no phenomeno da \emph{refracção}, que se dá quando os raios luminosos, passando de um para outro meio, mudam de direcção.\par Expliquemos a 2.ª e a 3.ª condição, relativas ao \emph{orgão} corporeo ou sensitivo. --- A 2.ª condição é que o nervo sensitivo tenha uma sufficiente excitabilidade, produzida pelo \emph{agente} nerveo. O \emph{agente nerveo} é uma qualidade activa, uma energia vital, que percorre os nervos á \emph{semelhança} de uma corrente electrica. Nos nervos \emph{motores}, esta energia dirige-se dos centros aos musculos, e por isso a sua direcção é \emph{centrifuga}; nos nervos \emph{sensitivos}, dirige-se da peripheria aos centros, e por isso a sua direcção é \emph{centripeta}. O \emph{agente nerveo} é apenas semelhante á \emph{electricidade}, porque d'ella differe pelas seguintes razões: 1.ª) a \emph{electricidade} é quasi tão veloz como a luz, ao passo que o \emph{agente} nerveo não percorre ao segundo senão 40 ou 60 metros, conforme os nervos são \emph{motores} ou \emph{sensitivos}; --- 2.ª) a energia \emph{electrica} não é detida na sua passagem pela ligadura do nervo, emquanto o agente \emph{nerveo} é detido; --- 3.ª) a energia electrica \emph{diminue} de intensidade á medida que se adianta nos fios electricos, ao passo que a energia nervea \emph{augmenta} em proporção do seu percurso. --- A 3.ª condição é que a faculdade sensitiva, para que possa receber em si a especie ou semelhança do objecto sensivel --- não deve ter de \emph{modo algum} as qualidades do mesmo objecto, --- ou não as deve ter \emph{no mesmo grau}. Esta condição é uma consequencia dos principios\pgn{267}expostos. Sendo o acto a expressão, a semelhança do agente, é da essencia do agente communicar a propria semelhança. Ora, como a semelhança se communica a quem a não tem, segue-se que um agente não pode actuar sobre um paciente perfeitamente semelhante. D'ahi o axioma: \emph{agens non agit in simile}. E nada mais rasoavel. Para se realizar uma acção, é preciso que o agente seja \emph{em acto}, e o paciente \emph{em potencia} com relação ao mesmo acto; porque, se o paciente estivesse \emph{em acto}, a acção do agente não tinha razão de ser. E' necessario, pois, que o agente e o paciente sejam dissemelhantes nalguma qualidade (que um dá e outro recebe), --- ou que, se ambos possuem a mesma qualidade, não a possuam no mesmo grau. D'ahi se segue que os sentidos, estando \emph{em potencia} para os sensiveis \emph{proprios}, não podem ter e não têem as qualidades \emph{physicas} de nenhum d'esses sensiveis. Assim a pupilla, que está \emph{em potencia} com relação ás côres, que é capaz de receber as côres, deve ser destituida de toda a côr; porque, se tivesse alguma côr, esta não deixaria ver nenhum outra côr. Concluimos que o sentido, para que possa receber as qualidades sensiveis segundo a \emph{forma intencional}, --- não as deve ter no ser \emph{physico}, que lhes é proprio. Não as deve ter \emph{de modo algum}, se o sentido é o da \emph{vista}, ou do \emph{ouvido}, ou do \emph{olfacto}, ou do \emph{gosto}; não as deve ter no \emph{mesmo grau}, se o sentido é o do \emph{tacto} (assim se o grau de temperatura da minha mão é egual ao grau de temperatura da agua, o calor da agua não impressiona a minha mão, nem pode modificar o seu calor) (S. Th., \emph{De Anima}, II, 1. 23).\par Dissemos nos n.os 155 e 157, e repetimos agora que, no phenomeno da sensação, ha a especie \emph{impressa} e a especie \emph{expressa}. Depois do que deixamos dito percebemos mais facilmente a \emph{natureza} de cada uma d'estas especies e a \emph{differença} entre uma e outra. Na verdade: 1.º) A especie \emph{impressa} é o effeito da acção do agente no paciente; por isso é a imagem d'essa acção, e é o principio immediato da percepção sensitiva. A especie \emph{expressa} é a reproducção vital da especie \emph{impressa}, feita pelo paciente; por isso é o resultado da reacção do paciente no agente, e é o \emph{termo} da percepção, não emquanto é o objecto da tendencia da faculdade, mas emquanto com elle acaba e nelle se completa o acto perceptivo. --- 2.º) A especie \emph{impressa}, sendo imagem do acto do agente e por isso sendo imagem do proprio agente, unida ao paciente, torna-o semelhante ao agente na ordem \emph{intencional}. A especie \emph{expressa}, sendo uma acção ou reacção vital do paciente, embora exprima ou represente o objecto, todavia exprime-o e representa-o de um modo conforme á indole ou natureza do paciente, e por isso torna-o semelhante ao proprio paciente; porque a \emph{especie}, por ser uma qualidade accidental, participa necessariamente da natureza do sujeito, a que adhere, e o objecto, por ella representado, é representado de um modo material ou espiritual, conforme material ou espirituel é o sujeito. --- Esta conclusão que o paciente, pela especie \emph{expressa}, torna semelhante a si o objecto --- não emquanto á \emph{representação}, --- mas emquanto ao \emph{modo} d'essa representação, é uma consequencia do grande principio: que todo o agente imprime no effeito a propria semelhança: \emph{omne agens agit sibi simile}. --- 3.º) A especie \emph{impressa} desapparece em pouco tempo. A especie \emph{expressa} fica na imaginação ou na\pgn{268}memoria num estado habitual, como representação viva do objecto percebido, e pode reviver por uma excitação mechanica dos nervos ou por um esforço da imaginação ou da memoria.\par Respondemos a algumas difficuldades.\par a) Dizem: Se a sensação é um phenomeno complexo, que resulta da impressão e da percepção, --- se ella é um acto unico, que tem dois aspectos differentes: a \emph{paixão} e a \emph{acção}, nada mais facil que se confunda a externo com o interno, o objectivo com o subjectivo.\par Esta difficuldade, que é a dos \emph{Cartesianos}, não tem razão de ser. Ainda que a acção e a paixão sejam no paciente um só acto, comtudo este acto conserva um duplice caracter --- um \emph{interno} e \emph{affectivo}, outro \emph{externo} e \emph{representativo}. E nós distinguimos perfeitamente o que é do sujeito e o que é do objecto, e não attribuimos a um o que é proprio do outro. Quando a ponta de uma agulha fere a minha mão, sei distinguir perfeitamente a ponta e a dôr, porque attribuo a ponta á agulha e a dôr a mim mesmo. E estas percepções são e se conhecem tão distinctas, que estão entre si na razão \emph{inversa}: quanto mais forte é uma, tanto mais fraca é a outra. Assim, uma luz muito viva cança os olhos e não deixa distinguir os objectos; ao passo que uma luz moderada torna mais perfeito o acto da visão. O Angelico Doutor (\emph{Quodl}. VII, a. 3) indica o criterio para a distincção dos dois elementos da sensação nas seguintes palavras: « Sensus exteriores \emph{suscipiunt tantum} (speciem) a rebus per \emph{modum patiendi}, sine hoc quod aliquid cooperentur ad sui formationem (phase \emph{passiva}); quamvis \emph{jam formati} (h. e. informati a specie ab objecto recepta) habeant \emph{propriam operationem}, quae est judicium de propriis objectis » (phase \emph{activa}).\par \emph{b}) Pode insistir-se: Os proprios escholasticos dizem que a percepção é \emph{relativa} á faculdade, e defendem o celebre axioma: tudo o que é recebido no paciente é recebido de um modo proporcionado ao mesmo paciente --- \emph{quidquid recipitur per modum recipientis recipitur}. Ora tudo isto quer dizer que o objectivo se confunde com o subjectivo, ou que o « \emph{objecto é como que transubstanciado na propria percepção} » (Mallet, \emph{Revue de Philos}., 1906, p. 242).\par A esta difficuldade, muito complexa e confusa, procuraremos dar uma resposta muito clara e convincente. --- Antes de tudo, admittimos a \emph{relatividade} na percepção. Por quanto, a existencia e a perfeição das percepções depende da existencia e da perfeição dos orgãos correspondentes. Quem não possue o orgão visual, ou o não possue em condições normaes, não pode perceber, ou não pode perceber perfeitamente os objectos coloridos. A percepção, portanto, é \emph{relativa} ao orgão correspondente. --- Mas do facto de a percepção ser imperfeita pela imperfeição do orgão, não se segue que ella deixa de ser \emph{objectiva}. O pouco que é percebido é realmente percebido, e existe realmente nas coisas. A percepção do mundo exterior é muito limitada, muito inadequada; mas deverá d'ahi concluir-se que não conhecemos de modo algum o proprio mundo? Não, por certo. --- Tambem concedemos e admittimos --- que a percepção depende, em grande parte, das disposições actuaes dos orgãos sensitivos, --- e que os varios graus de desattenção, de cançaço, de exgottamento, de impotencia derivada das doenças ou da edade, enfraquecem\pgn{269}a percepção dos sentidos. Esta \emph{relatividade} é muito natural e é uma consequencia da theoria da \emph{acção} e da \emph{paixão}. O paciente não pode receber a acção do agente senão na medida da sua capacidade e das suas disposições. Por isso é necessaria uma certa proporção entre o agente e o paciente; como tambem é necessario remover todo e qualquer obstaculo, que possa impedir a acção de um no outro (aliás a acção ou não será recebida, ou chegará desfigurada). Quando esta proporção existe, e a sensação se realiza em condições normaes, não só não confundimos o objectivo com o subjectivo, mas distinguimos facilmente um elemento do outro, porque a \emph{acção do agente não muda de natureza}, \emph{quando é recebida no paciente}. --- E já temos a resposta á ultima parte da difficuldade. Diz-se um absurdo, quando se diz que a acção do agente muda de natureza, quando e porque é recebida no paciente; e este absurdo nasce do facto de confundirem os dois estados do movel ou paciente, que são a \emph{paixão} e a \emph{reacção}. O acto do agente, diz S. Thomaz, \emph{não se torna conforme} com a \emph{natureza} do paciente, mas é apenas \emph{proporcionado á capacidade} do paciente. Por isso, assim como a agua não muda de natureza pelo facto de ser recebida numa garrafa, embora não possa ser recebida senão numa quantidade proporcionada á capacidade da garrafa, assim tambem não muda de natureza, não se \emph{transubstancia} no paciente o acto do agente pelo facto de ser recebido no paciente, embora deva accomodar-se ás condições do proprio paciente. Eis as palavras do S. Doutor: « Cum dicitur \emph{unumquodque esse in alio secundum modum ejus in quo est}, \emph{intelligitur quantum ad capacitatis ipsius modum}, \emph{non quantum ad naturam ejus}. Non enim oportet ut id quod est in aliquo habeat \emph{naturam} et \emph{proprietatem} ejus in quo est, sed quod recipitur in eo secundum \emph{capacitatem} ipsius. Manifestum est enim quod aqua non habet naturam amphorae » (\emph{De Anima}, q. 1, a. 10 ad 14). --- Esta doutrina é exposta pelo Angelico em muitas partes das suas obras, dizendo que os objectos conhecidos se encontram na cognoscente de um modo proporcionado á natureza d'este, e por isso as coisas \emph{materiaes} são recebidas na intelligencia de um modo \emph{espiritual}, isto é, segundo a capacidade da propria intelligencia, que é faculdade \emph{espiritual}.\par \emph{c}) Alguns physiologistas dizem que a percepção \emph{sensitiva} não passa de uma \emph{impressão organica}, que se torna \emph{consciente}; e os materialistas, antigos e modernos, reduzem-na a um mero \emph{movimento mechanico}.\par A opinião d'esses physiologistas é falsa. Se a percepção \emph{sensitiva} fosse apenas uma \emph{impressão organica}, seria um phenomeno meramente \emph{passivo}. Ora isto, como vimos, é falso; porque o sujeito sensiente não é meramente passivo, mas reage e completa a acção dos objectos externos, tornando-os semelhantes a si mesmo, emquanto ao \emph{modo} por que os percebe. Nem vale dizer que a impressão organica se torna \emph{consciente}; porque dão-se muitas percepções, que se não tornam conscientes e que, todavia, são verdadeiras percepções. --- Não insistimos na refutação do erro dos materialistas, pois foi refutado neste artigo e na \emph{Cosmologia especial}, quando demonstrámos que os phenomenos \emph{vitaes} (e a sensação é um phenomeno \emph{vital}) não são meros movimentos \emph{mechanicos}. O movimento \emph{mechanico} é effeito de um principio \emph{exterior}; ao passo que o movimento \emph{vital}, embora possa ser determinado por\pgn{270}um impulso exterior, é uma operação \emph{immanente}, que tem o seu principio e o seu termo no mesmo ser vivo. Não negamos que as percepções sensitivas são acompanhadas de vibrações nervosas, de movimentos mechanicos; mas é absurdo identificar dois entes, pelo facto de andarem acompanhados.\par Devemos advertir, para evitar equivocos, que os escriptores de coisas Philosophicas não concordam na significação da palavra « \emph{sensação} ». Os antigos, empregando esta palavra, queriam denotar \emph{ambos} os elementos do phenomeno \emph{perceptivo-sensitivo}: --- o elemento \emph{passivo}, que é a \emph{impressão organica} ou a \emph{modificação} do sujeito sensiente, e o elemento \emph{activo}, que é a propria \emph{percepção}. Muitos modernos, não querendo seguir a tradição dos antigos, reservam a palavra « \emph{sensação} » para indicar exclusivamente os actos do \emph{appetite sensitivo}. --- A questão é de palavras, e resolve-se declarando o sentido, em que o vocabolo se emprega. Para maior clareza de exposição, empregamos a palavra « \emph{impressão} » para exprimir o elemento \emph{passivo} do phenomeno \emph{perceptivo-sensitivo}, e a palavra « \emph{percepção} » para denotar o seu elemento \emph{activo}, e, algumas vezes, quando não havia perigo de equivocos, servimo-nos da palavra « \emph{sensação} » para indicar todo o phenomeno.}.

b) Propriedades da percepção sensitiva externa. --- As principaes propriedades da percepção \emph{sensitiva externa} são: a \emph{qualidade}, a \emph{quantidade}, a \emph{duração}, e a \emph{tonalidade}. --- A \emph{qualidade} da \pg{267}  sensação depende da qualidade do estimulo externo, ou do objecto agente, porque os actos especificam-se pelos objectos; assim temos o phenomeno da \emph{visão}, da \emph{audição}, etc., conforme \pg{268}  a \emph{luz} excita a \emph{vista}, o \emph{som} o \emph{ouvido}, etc. --- A \emph{quantidade}, ou \emph{intensidade}, da sensação depende do grau, maior ou menor, da energia, com que a faculdade perceptiva experimenta a impressão \pg{269} , proveniente do estimulo externo. --- A \emph{duração} da sensação depende da duração do estimulo externo; embora, geralmente, o principio e o termo do estimulo preceda, de

algum tempo, o principio e o termo da sensação. --- A \emph{tonalidade} é a modificação de prazer ou de dôr, que a sensação produz no sujeito sensiente\footnote{Deixando de fallar na \emph{qualidade} da sensação, pois a qualidade da sensação depende, como dissemos, da qualidade dos objectos, e já vimos \emph{quaes} são os objectos que impressionam ou modificam os nossos sentidos externos, --- e reservando para outro artigo a questão ácerca da \emph{tonalidade} da mesma sensação, tratamos agora da sua \emph{intensidade} e \emph{duração}.\par a) A \emph{intensidade} da sensação é o grau maior ou menor de força, com que a faculdade respectiva prova ou experimenta a sensação.\par Poderá essa \emph{intensidade} ser medida?\par Está averiguado que existe uma \emph{certa} dependencia entre a intensidade da sensação o a intensidade da excitação que a provoca; assim a luz de duas velas é mais intensa que a de uma só. Mas também é certo que a sensação não está na razão \emph{directa} da excitação; assim, num salão illuminado por mil velas, a sensação da luz não é mais forte, accendendo uma outra vela. Qual será, pois, a relação entre o augmento ou diminuição da excitação e a maior ou menor intensidade da sensação? Ou, por outras palavras: que quantidade se deve accrescentar á causa da sensação, ou tirar d'ella, para que se perceba a differença, para mais ou para menos, na propria sensação?\par E' este o objecto da \emph{psyco-physica}.\par E' certo que a quantidade, que se deve accrescentar á causa excitante, para provocar uma modificação perceptivel na sensação, não é \emph{absoluta}, mas \emph{relativa} á quantidade da propria causa.\par \emph{Weber} diz que o augmento da excitação, o qual deve produzir uma modificação perceptivel no phenomeno sensitivo, está numa \emph{relação constante} com a quantidade da causa excitante, a que elle se accrescenta: --- ou, por outras palavras: existe uma \emph{relação constante} entre o peso \emph{inicial}, que provoca uma sensação, e o peso \emph{addicional}, que deve causar um augmento de intensidade na mesma sensação. Essa \emph{relação constante} é expressa por ¹/₃. Assim, para a differença de um peso 1 se tornar perceptivel, é necessario\pgn{271}acrescentar um peso ¹/₃, a um peso 2 é necessario acrescentar um peso ²/₃, a um peso 3 é necessario acrescentar um peso 1, etc.\par \emph{Fechner}, querendo generalisar o calculo de \emph{Weber} e elevál-o a uma lei mathematica, disse que, para a sensação augmentar em progressão \emph{arithmetica}, é necessario que a excitação augmente em progressão \emph{geometrica}.\par As conclusões de \emph{Weber} e de \emph{Fechner} foram diversamente apreciadas pelos sabios. Uns, como \emph{Hering} e \emph{Delboeuf}, emendaram-nas; --- outros rejeitaram-nas, como destituidas de todo o fundamento, pois que (segundo estes) o problema é insoluvel; --- e outros, como \emph{Wundt}, \emph{Helmholtz}, \emph{Ebbinghaus}, admittindo como hypothese a lei de \emph{Weber}, procuraram novos methodos, para a solução do problema que forma o objecto da \emph{psyco-physica}.\par Posto isto, façamos duas perguntas:\par 1.º) Será insoluvel o problema ácerca da \emph{medida das sensações}?\par 2.º) E, se não for insoluvel, que deve dizer-se da lei de \emph{Weber}?\par Resposta á 1.ª pergunta:\par Não podemos medir \emph{exactamente} as sensações. Porquanto, para isso, deveria encontrar-se uma \emph{unidade} de medida, a que se podessem reduzir todas as sensações. Mas esta \emph{unidade} de medida não existe --- nem para as sensações de differente \emph{qualidade} (pois \emph{ver}, \emph{gostar}, \emph{ouvir}, etc. são phenomenos irreductiveis \emph{a} uma medida commum), --- nem para as sensações da mesma \emph{qualidade}, mas de differente \emph{intensidade} (pois, embora percebamos que uma sensação é mais forte e mais intensa que outra, não podemos dizer \emph{quantas vezes} aquella é mais forte e mais intensa que esta).\par Mas, se não podemos medir exactamente as sensações, podemos comtudo chegar a determinar, mais ou menos approximadamente, qual deve ser o augmento da excitação na causa, para haver um augmento na sensação. Com effeito, experimentando uma certa e determinada sensação, se augmentarmos a quantidade na causa que a excita, perceberemos uma differença na propria sensação. Marcando a quantidade, que foi necessaria para a variação da sensação, e repetindo a experiencia, teremos novamente a quantidade, que foi necessaria para a modificação perceptivel da mesma sensação. Renovem-se muitas vezes estas experiencias, mudando as circumstancias de tempo, de pessoas e de meio, a \emph{media} de todos os differentes resultados indicará a quantidade, que é necessario accrescentar a uma causa excitante, para haver uma modificação \emph{perceptivel} na sensação.\par Resposta á 2.ª pergunta.\par A universalidade da lei de \emph{Weber} não está, até hoje, confirmada pela experiencia. As comparações feitas entre as diversas sensações e a lei de \emph{Weber} dão os seguintes resultados. As sensações \emph{auditivas} são as que mais concordam com essa lei; --- o accordo entre as sensações de pressão, de movimento e de gosto e a mesma lei é mais limitado; --- as sensações \emph{thermicas}\pgn{272}dão resultados mais incertos. --- As sensações \emph{olfactivas} e \emph{geraes} não se prestam ás experiencias. --- Mas, embora a universalidade da lei de \emph{Weber} não esteja até hoje confirmada pela experiencia, e embora seja difficil tal confirmação relativamente a todas as sensações, comtudo não se lhe pode recusar, em \emph{nome da experiencia}, o caracter de lei \emph{geral}.\par b) A \emph{duração} da sensação é um outro problema, que está sendo vivamente discutido pelos sabios, mas que está longe de ser resolvido.\par Quanto tempo emprega um phenomeno \emph{physiologico} para percorrer o \emph{arco reflexo}, que elle descreve?\par Mediu-se esse tempo com relação ás sensações da \emph{vista}, do \emph{ouvido}, do \emph{tacto}, e também do \emph{gosto} e do \emph{olfacto}.\par Eis o resultado das experiencias feitas.\par A sensações da \emph{vista}, para \emph{Donders}, duram ¹/₅ de segundo, --- e, para \emph{Beaunis}, ²⁵/₁₀₀ de segundo.\par A sensações do \emph{ouvido}, na opinião de \emph{Donders}, duram ¹/₆ de segundo, --- e, na opinião de \emph{Beaunis}, 15.9 centesimos de segundo.\par As sensações do \emph{tacto}, segundo \emph{Donders}, gastam ¹/₇ de segundo para percorrer o \emph{arco reflexo}, --- e, segundo \emph{Beaunis}, 10.6 centesimos de segundo.\par Na sensações do \emph{gosto} e do \emph{olfacto}, o tempo \emph{physiologico} ou de \emph{reacção} varia consideravelmente.}.

%%>>> t160

\pg{266} \pg{270} \pg{271} \tese{160}{Séde da percepção sensitiva externa.}
A séde da percepção \emph{sensitiva externa} não é o cerebro, mas é o proprio \emph{orgão animado externo}. Por quanto, a experiencia interna, confirmada pelo senso commum, attesta-nos --- que vemos as \emph{côres} com os \emph{olhos}, e não com o cerebro, --- que ouvimos os \emph{sons} com os \emph{ouvidos}, e não com o cerebro, --- que sentimos a resistencia dos corpos com a \emph{cutis} e as \emph{mucosas}, e não com o cerebro, etc.

--- Esta convicção encontra tambem apoio na \emph{anatomia}, a qual demonstra que cada um dos sentidos externos é dotado de um apparelho nervoso especial, natural e essencialmente adaptado para ser o unico sujeito receptivo e perceptivo de cada um dos sensiveis \emph{proprios}; assim o apparelho \emph{visual}, e só elle, é capaz de receber e perceber as \emph{côres}, --- o apparelho \emph{auditivo}, e só elle, é capaz de receber e perceber os \emph{sons}, etc. Logo a séde da percepção sensitiva externa não é o cerebro, mas é o proprio orgão animado externo \footnote{Alguns physiologistas e philosophos sustentam que todas as sensações \emph{externas} se realizam no \emph{cerebro}, e que os orgãos \emph{externos} são apenas necessarios para receber as impressões, que os nervos meramente \emph{conductores} transmittem ao cerebro. E' o \emph{encephalismo}, isto é, o systema que reconhece o cerebro como unico orgão de todas as sensações externas, de modo que é o cerebro que vê, que ouve, etc. --- Esta opinião é uma das varias consequencias da theoria \emph{cartesiana}, que colloca a alma numa cellula central do cerebro e sustenta que todos os outros nervos têm apenas o officio de fios telegraphicos ou de orgãos de transmissão.\par Mas a theoria é tão falsa, como é falso o fundamento, em que se basêa. --- Não só não negamos, mas até defendemos que o cerebro deve perceber e percebe de algum modo as sensações \emph{externas}, porque é elle a séde das sensações \emph{internas}, que dependem das \emph{externas}, e porque é orgão central, de que, como de uma \emph{raiz commum}, deriva, e para o qual, como para o seu termo, converge toda a energia nervea, informante os sentidos externos. Diz S. Thomaz: « Vis sentiendi diffunditur in organa quinque sensuum ab aliqua \emph{una radice communi}, a qua procedit vis sentiendi in omnia organa, ad quam etiam terminantur omnes immutationes singulorum organorum » (\emph{De Anima}, III, l. 3). --- Admittimos tambem que as \emph{sensações externas}\pgn{273}\emph{conscientes}, ou \emph{advertidas}, dependem, de algum modo, do cerebro, --- não emquanto são \emph{percepções externas}, mas emquanto são \emph{conscientes}; porque o cerebro, como vimos, é o orgão da consciencia sensitiva e da memoria sensitiva, e os nervos sensitivos communicam com os emispherios cerebraes. --- Mas não admittimos a opinião que colloca a séde das sensações \emph{externas} no cerebro, e não nos sentidos externos. O \emph{senso commum}, que se funda na evidencia, e a experiencia \emph{interna} revoltam-se contra semelhante theoria. Na verdade, se o cerebro fosse a séde de todas as sensações \emph{externas}, elle deveria possuir um apparelho adaptado para perceber as \emph{côres}, e outro para os \emph{sons}, e outro para os \emph{sabores}; porque sensiveis especificamente diversos não podem ser percebidos por uma só faculdade. Ora o cerebro não tem esses apparelhos. --- Alem d'isso, se o cerebro fosse a séde das sensações \emph{externas}, é evidente que um animal, a que fosse arrancado o cerebro, não poderia ter nenhuma sensação \emph{externa}. Ora os modernos physiologistas sustentam que a ablação do cerebro não é sufficiente para destruir completamente o exercicio dos sentidos em geral, e da \emph{vista} e do \emph{ouvido} em particular, porque averiguaram que os animaes, aos quaes tinham arrancado os emispherios, evitavam obstaculos, faziam movimentos repentinos na occasião de fortes rumores, seguiam com a cabeça o movimento de uma vela accesa, davam signaes de que viam e ouviam. --- Repetimos que uma coisa é a percepção \emph{externa}, outra é a \emph{consciencia} d'essa percepção; esta depende do cerebro, não aquella. --- D'ahi a divisão, introduzida por \emph{Munk} e seguida por \emph{Luciani}, da \emph{cegueira}, como tambem da \emph{surdez}, em \emph{externa} e \emph{psychica}. A \emph{cegueira} e a \emph{surdez psychica} consistem em que o animal vê ainda as côres e ouve os sons, mas não adverte que vê e que ouve. --- D'ahi tambem a distincção das sensações em \emph{inconscientes}, e \emph{conscientes}. Estas dependeriam somente dos emispherios cerebraes; aquellas dependeriam do mesencephalo, da ponte e da medulla allongada e espinhal. (Cf. Luciani e Serpilli, \emph{Le localizzazioni funzionali del cervello}). --- De resto, o proprio \emph{Richet}, defensor acerrimo do \emph{nervo meramente conductor}, confessa que a sua theoria « é um dogma sem prova, e um dia pode tornar-se uma idea falsa » (\emph{Physiol}. \emph{des nerfs}).\par Os adversarios fazem diversas difficuldades. Examinemos algumas.\par a) Dizem: A sensação não pode realizar-se senão nas cellulas de substancia \emph{cinzenta}, e esta só se encontra no cerebro. --- Respondemos que estas affirmações são falsas. Os melhores physiologistas reconhecem que ha cellulas de substancia \emph{cinzenta} tambem nos nervos, sobretudo nas extremidades d'elles, as quaes se encontram nos orgãos dos sentidos externos (Cf. Luys, \emph{Le cerveau}). Além d'isso, todos admittem que a substancia \emph{cinzenta} do cerebro é \emph{insensivel}, e que, se a impressão chega ao cerebro por um meio distincto dos nervos, não se realiza a percepção.\par b) Continuam: Está averiguado que os nervos sensitivos locaes devem estar em communicação com o \emph{cerebro}, para que possa realizar-se á percepção. Ora isto mostra que é propriamente o cerebro que sente, e que os nervos sensitivos locaes têm apenas o officio de transmittir ao cerebro as impressões externas. --- Respondemos ser verdade que, para a percepção sensitiva, é indispensavel que os nervos sensitivos locaes communiquem com os \emph{centros nervosos}. Mas este phenomeno mostra --- não que são os \pgn{274}centros nervosos que directamente percebem as qualidades sensiveis dos corpos, --- mas que os nervos sensitivos se tornam paralyticos, isto é, perdem toda a excitabilidade necessaria para a percepção sensitiva, se forem separados dos \emph{centros nervosos}, que são a fonte da actividade nervosa. Dissemos que, para a percepção sensitiva, devem os nervos sensitivos communicar com os \emph{centros nervosos}, e não com \emph{o cerebro}; porque, como vimos, a ablação dos \emph{emispherios cerebraes} não é sufficiente para destruir inteiramente o exercicio dos sentidos externos.\par c) Insistem: Um homem, a quem foi cortado um braço, continua a localisar a dor no braço, que já não tem. Ora do mesmo modo nós attribuimos a sensação ás extremidades nerveas dos orgãos externos, embora essa sensação se realize exclusivamente no cerebro. --- Respondemos, negando a paridade dos casos. No primeiro caso, ha uma illusão de imaginação; no segundo, não ha. Na verdade, o homem, a quem foi cortado o braço, não sente realmente a dor no braço cortado, mas sim na parte, que antes estava unida á que foi cortada, e, se elle attribue, ainda per algum tempo, á dor á parte cortada, isso depende --- não só do habito que tinha de sentir naquelle modo e de perceber o seu braço inteiro, --- mas tambem da associação das imagens sensitivas das proprias sensações, que uma vez se experimentavam na parte cortada. Pelo contrario, a convicção commum de que as sensações externas se realizam nos proprios orgãos externos, não se explica pelo habito, nem pela associação dos phantasmas, --- e não cessa com o andar do tempo; o que prova que não é um engano da imaginação, mas que se funda na evidencia objectiva das coisas.\par *\par Depois de termos tratado da percepção dos sentidos \emph{externos} em geral, deveriamos agora tratar da percepção de cada sentido \emph{externo} em particular. Mas, tendo já indicado o \emph{orgão} de cada sentido \emph{externo}, o seu \emph{objecto} e \emph{mechanismo}, não nos detemos mais neste argumento. --- Só queremos tocar num ponto, que é muito contestado e que se refereá \emph{necessidade} da \emph{especie sensivel} para \emph{cada um} dos sentidos \emph{externos}. Uns affirmam que a \emph{especie sensivel} é necessaria para \emph{cada um} dos sentidos \emph{externos}; outros sustentam que só é necessaria para os sentidos da \emph{vista} e talvez do \emph{tacto}, porque os \emph{sensiveis proprios} dos outros sentidos \emph{externos} podem unir-se com estes na sua physica realidade. Seguimos a \emph{primeira} opinião, que defende a \emph{necessidade} da \emph{especie sensivel} para \emph{todos} e \emph{cada um} dos sentidos \emph{externos}. Na verdade, os \emph{sensiveis proprios} são todos \emph{objectivos}, isto é, são qualidades \emph{activas}, que existem nos proprios objectos externos. Sendo \emph{objectivos}, os \emph{sensiveis proprios}, para serem percebidos, devem actuar nos respectivos sentidos externos e reduzil-os da potencia ao acto. Ora os sentidos \emph{externos}, sendo por si indeterminados, não podem ser reduzidos da potencia ao acto, se não forem completados e determinados ao proprio acto, e não podem ser completados e determinados senão por uma especie ou forma; porque, como dissemos, é sempre a forma, substancial ou accidental, que completa a potencia e á determina a uma certa especie ou operação. --- Este razão vale --- não só para a \emph{vista}, mas para \emph{todos} os sentidos, porque todos são indeterminados, e por isso todos precisam de uma forma, ou especie, ou imagem, que os determine a perceber este ou aquelle objecto. --- Nem se diga que a determinação deriva da proprio objecto externo, emquanto actua na faculdade. O objecto externo não pode ser forma complementar e determinativa do sentido pela sua \emph{physica realidade}, mas pode sel-o por uma sua \emph{forma}, ou \emph{especie}, ou \emph{imagem}, que elle produz e imprime na faculdade sensitiva. --- Nas notas ao n.º 156 dissemos, com S. Thomaz, --- que a impressão do objecto nos sentidos produz uma duplice mudança: uma \emph{natural} e outra \emph{intencional}, --- que na mudança \emph{intencional} a forma ou especie do agente é recebida no sentido, --- e que esta especie é indispensavel para o conhecimento sensitivo.\par Tendo provado, em geral, a \emph{necessidade} da especie \emph{sensivel} para \emph{todos} os sentidos \emph{externos}, agora digamos alguma coisa ácerca da \emph{especie sensivel} de \emph{cada} sentido \emph{externo}. --- Em quanto ao \emph{tacto}, as especies são necessarias, porque o \emph{orgão} d'este sentido não é propriamente a \emph{epiderme}, isto é, a superficie exterior da cutis, mas é a \emph{derme}, isto é, o tecido subjacente á epiderme; e por isso o objecto externo, não actuando immediatamente no sentido pela sua realidade physica, deve actuar nelle por uma sua imagem ou especie. Chega-se á mesma conclusão, se se fixar, com alguns physiologistas, a séde d'este sentido nas \emph{fibras nervosas}; porque, tambem nesta hypothese, o objecto não actua immediatamente senão nos musculos ou noutra parte externa ou interna. --- Em quanto ao \emph{gosto}, alguns negam a necessidade da \emph{especie sensivel}, porque o objecto une-se immediatamente, por si mesmo, ao sentido, exercendo nelle uma acção de ordem chimica. Mas, se o objecto do gosto se une, immediatamente e na sua realidade physica, com o \emph{orgão}, não se une com o \emph{sentido}; com este deve unir-se pela sua \emph{imagem}. --- Diga-se o mesmo emquanto ao \emph{olfacto}. --- Em quanto ao \emph{ouvido}, o seu objecto é o som. O som é uma qualidade activa dos corpos, sempre em potencia e ás vezes em acto, e resulta das vibrações dos corpos e propaga-se por um meio, que, segundo alguns, é o ar, e, segundo outros, é a agua, ou um corpo aeriforme, ou um corpo liquido ou solido. O \emph{sujeito}, que emitte o som, não é o ar, nem é outro qualquer meio transmissor, mas é o proprio corpo. Sendo o som uma qualidade do corpo sonoro, deve unir-se com os ouvidos por uma especie ou imagem \emph{intencional}. --- Em quanto á \emph{vista}, o seu objecto é constituido pela \emph{luz} e pela \emph{côr}, que são qualidades objectivas dos corpos. Uns dizem que a côr se distingue da \emph{luz}; mas outros affirmam que a côr não se distingue da \emph{luz}, mas que é apenas uma modificação da luz, isto é, que a côr é a propria luz diversamente modificada, segundo a diversidade da superficie do corpo opaco. --- Será a \emph{côr} uma qualidade \emph{permanente} dos corpos? Os que distinguem a côr da luz, dizem que sim, embora accrescentem que ella --- a côr --- não pode ser vista sem a luz. Os que identificam a luz com a côr, affirmam que nos objectos --- não existe a côr, tomada no seu sentido \emph{formal}, emquanto resulta do actual reflexo de certos raios, porque, então, não ha côr sem luz, --- mas existe a côr \emph{fundamental}, emquanto a natureza ou estructura do corpo exige uma certa e determinada côr pela absorpção de certos raios e pela reflexão dos outros, e esta côr existe permanentemente nos corpos, mesmo na ausencia da luz. --- Mas a luz e a côr, quer se distingam entre si, quer não, constituem sempre um só objecto da \emph{vista}. --- Ora para \pgn{276}o acto da \emph{visão} são indispensaveis as \emph{especies} ou \emph{imagens sensiveis}. Na verdade, o objecto da \emph{vista}, que é proprio ente externo, ou a luz, ou a côr, não pode unir-se immediatamente por si com esse sentido, mas limita-se a impressionar a retina, e por isso deve unir-se com a mesma faculdade por meio de uma sua imagem \emph{intencional}.}.

%%>>> t161

\tese{161}{Corollario.}
Se a séde da percepção sensitiva externa é o proprio sentido \emph{externo}, e não o cerebro, é evidente que \pg{272} \pg{273} \pg{274}  podem dar-se sensações \emph{inconscientes}, isto é, sensações de que não temos consciencia. --- Para que haja verdadeira \emph{sensação}, \pg{275}  basta o respectivo sentido externo; para que a sensação seja \emph{consciente}, é necessario o cerebro.

%%>>> art-23

\section*{Artigo III.\\ Percepção sensitiva interna.}
\addcontentsline{toc}{section}{Artigo III. --- Percepção sensitiva interna.}

%%>>> t162

\pg{276} \tese{162}{Percepção sensitiva interna.}
A percepção \emph{sensitiva interna} é o acto proprio dos sentidos \emph{internos}, pelo qual o homem conhece, de um modo \emph{material}, ou \emph{concreto}, seus actos e estados \emph{sensitivos} e \emph{internos} \footnote{No artigo II do capitulo II da secção II d'este tratado, fallámos, com o conveniente desenvolvimento, dos sentidos \emph{internos} e dos seus respectivos \emph{actos}. Limitamo-nos, portanto, n'este artigo, a indicar o modo, por que esses sentidos operam.}.

%%>>> t163

\tese{163}{Objecto da percepção sensitiva interna.}
O \emph{objecto} da percepção \emph{sensitiva interna} é constituido pelos nossos \emph{actos} e \emph{estados sensitivos} e \emph{internos}. Esses \emph{actos} e \emph{estados} formam o objecto \emph{immediato} ou \emph{explicito}; porque o objecto \emph{mediato} ou \emph{implicito} é constituido pelos \emph{agentes externos} ou pelo proprio \emph{sujeito sensiente} \footnote{Os sentidos \emph{internos} não soffrem a acção dos agentes \emph{externos}, como estes são em si mesmos, mas como são recebidos e existem nos sentidos \emph{externos}, e por isso sem materia. O acto, portanto, dos sentidos \emph{internos} não se refere \emph{immediatamente} aos \emph{objectos} ou \emph{agentes externos}, mas ao \emph{acto} dos sentidos externos (S. Th. \emph{De Anima}, III, l. 2). Com effeito, o homem, pela \emph{imaginação} ou pelo \emph{sentido commum}, --- não vê as côres, nem ouve os sons, --- mas \emph{sente} ou \emph{percebe} que vê as côres e que ouve os sons, e emtanto chega á percepção dos objectos externos, emquanto o ver e o ouvir são apenas uma reproducção vital dos proprios objectos, que actuaram nos sentidos da vista e do ouvido. Por isso não pode receber-se a \emph{imagem} da sensação, sem que se conheça, ao mesmo tempo, o seu \emph{objecto}. --- Veja-se no logar indicado o que dissemos ácerca do \emph{objecto} de cada sentido \emph{interno}.}.

%%>>> t164

\tese{164}{Os sentidos internos precisam de uma imagem para o seu acto.}
Os sentidos \emph{internos}, como os \emph{externos}, são por si indifferentes ou indeterminados com relação aos muitos e vários objectos externos, que podem perceber, e d'essa indifferença ou indeterminação só podem sahir pela acção dos proprios objectos. Ora os objectos externos, se não podem actuar nas faculdades \emph{internas} na sua physica realidade, devem actuar por meio de uma imagem, que os represente. Logo os sentidos \emph{internos} precisam de uma imagem para o seu acto \footnote{A conclusão é evidente. Se os sentidos \emph{externos} precisam de uma\pgn{277}imagem ou especie do objecto para sahir da sua indifferença ou indeterminação, muito mais precisam de uma imagem os sentidos \emph{internos}, que não só têm a mesma indifferença ou indeterminação, mas tambem se encontram no interior do organismo e escapam á acção directa e immediata dos objectos externos. --- Esta imagem ou especie é indispensavel --- não só quando o objecto está \emph{ausente}, mas tambem quando está \emph{presente}, isto é, quando \emph{actualmente} influe nos mesmos sentidos externos. Por quanto, a imagem do objecto \emph{presente}, se pode unir-se com os sentidos \emph{externos}, não se pode unir com os \emph{internos}. Cada faculdade deve ter a sua imagem ou especie proporcionada, que a complete em ordem ao acto; e por isso, sendo os sentidos \emph{internos} faculdades especificamente diversas dos sentidos \emph{externos}, as imagens ou especies, que informam os sentidos \emph{internos}, devem ser especificamente diversas das imagens ou especies, que informam os sentidos \emph{externos}. --- Essa differença manifesta-se no facto --- que, emquanto as especies dos sentidos \emph{externos}, dependem, na sua existencia, da presença physica do objecto externo e desapparecem ao desapparecer do proprio obiecto, as especies dos sentidos \emph{internos} não dependem, na sua existencia, da presença physica do objecto externo e, quando este desapparece, ellas não desapparecem. --- Nem se objecte que, para determinar ao acto os sentidos \emph{internos}, sobretudo a \emph{imaginação} e a \emph{memoria}, basta a presença da sensação \emph{externa}. A presença da sensação \emph{externa} não pode por si determinar e completar os sentidos \emph{internos}, porque essa sensação tem a sua séde no sentido \emph{externo}, e não pode unir-se, por si, com os sentidos \emph{internos}. E se, posta a sensação \emph{externa}, os sentidos \emph{internos} da \emph{imaginação} e da \emph{memoria} produzem o seu acto, este phenomeno se realiza --- não porque a sensação \emph{externa} determina e completa, por si, essas faculdades, --- mas porque, posta a sensação \emph{externa}, reproduzem-se imagens ou representações, que, unidas com a \emph{imaginação} e com a \emph{memoria}, completam estes sentidos e os reduzem ao acto.}.

%%>>> t165

\pg{277} \tese{165}{Essa imagem é produzida pela sensação externa.}
Attesta-nos a consciencia que, pelos sentidos \emph{internos}, percebemos a paixão ou modificação dos sentidos \emph{externos} e os seres que de algum modo actuaram nos mesmos sentidos. Ora a imagem, que reproduz a paixão ou modificação dos sentidos \emph{externos} e os seres que de algum modo actuaram nos mesmos sentidos, só pode ser produzida pela sensação \emph{externa}, porque só esta representa essa modificação e esses seres. Logo a imagem, de que os sentidos \emph{internos} precisam para o seu acto, é produzida pela sensação \emph{externa} \footnote{Os sentidos \emph{internos} precisam de uma especie ou imagem para o seu acto. Esta imagem não pode ser produzida --- pelo \emph{objecto externo}, --- nem pela \emph{imagem} dos sentidos \emph{externos}, --- mas só pela \emph{sensação externa}. --- Não pode ser produzida pelo \emph{objecto externo}, porque este não pode actuar immediata ou directamente nos sentidos \emph{internos}; nem pela \emph{imagem} ou \emph{especie} (\emph{impressa}) dos sentidos \emph{externos}, porque tal imagem, devendo ser especificamente diversa da imagem dos sentidos \emph{internos}, não pode completar estes sentidos; mas deve ser produzida pela \emph{sensação externa}, a qual, por ser uma qualidade activa, pode actuar nos sentidos \emph{internos} e imprimir nelles uma\pgn{278}especie ou imagem, que a represente, isto ó, que represente a modificação dos sentidos \emph{externos} e os objectos que a produziram. --- Esta acção da sensação \emph{externa} nos sentidos \emph{internos} é muito natural; porquanto os sentidos \emph{externos}, embora sejam distinctos dos \emph{internos}, estão todavia unidos com os \emph{internos} pelo systema nervoso, e por isso a sensação \emph{externa} pode e deve propagar-se até aos orgãos dos sentidos \emph{internos}, produzindo nestes uma especie ou imagem que a represente. E não repugna que a sensação \emph{externa} produza uma imagem, que a represente e que seja principio de uma outra percepção. Tambem o phantasma é uma especie ou imagem da coisa percebida, e todavia concorre para a formação da especie intelligivel, que é o principio do conhecimento intellectual. --- Nem se diga que a propria sensação \emph{externa} pode por si unir-se, como especie, com os sentidos \emph{internos}. A sensação \emph{externa}, como dissemos, é um acto vital que tem o seu principio e o seu termo no proprio sentido \emph{externo}, e por isso, se pode influir nos sentidos \emph{internos} e produzir nelles uma especie que a represente, não lhes se pode unir como especie. --- Como se vê, é sempre o objecto \emph{externo} que actua nos sentidos \emph{externos} e nos \emph{internos}, ainda que de um modo differente. Nos sentidos \emph{externos} o objecto actua \emph{directamente}, ao passo que nos \emph{internos} actua \emph{indirectamente}, isto é, por meio da sensação externa. A imagem, que o objecto externo imprime nos sentidos \emph{externos}, representa \emph{exclusiva} e \emph{directamente} o proprio objecto, porque a modificação ou o estado do sujeito é percebido por outra faculdade, que é o \emph{sentido commum}; mas a imagem, que o objecto externo produz, por meio da sensação \emph{externa}, nos sentidos \emph{internos}, \emph{directamente} representa a modificação ou o estado do sujeito, e \emph{indirectamente} representa o objecto \emph{externo}.}.

%%>>> t166

\pg{278} \tese{166}{Os sentidos internos, informados pela imagem, percebem o objecto.}
O objecto dos sentidos \emph{internos} é duplice: \emph{directo} e \emph{indirecto}. O \emph{directo} é a \emph{modificação} ou \emph{estado} do sujeito sensiente; o \emph{indirecto} é o \emph{objecto externo}, que provocou essa modificação ou esse estado. Ora diz-nos a experiencia que os sentidos \emph{internos}, logo que sejam actuados pela imagem, produzida pela sensação \emph{externa}, percebem, embora de um modo differente, não só a \emph{modificação} ou o \emph{estado} dos sentidos \emph{externos}, mas tambem o \emph{objecto externo}, que provocou essa modificação ou esse estado. Logo os sentidos \emph{internos}, informados pela imagem, percebem o objecto \footnote{A sensação \emph{externa} é \emph{subjectiva} e \emph{objectiva}. Por isso a imagem, produzida no sentido \emph{interno} pela sensação \emph{externa}, reproduz e exprime este duplice aspecto, isto é, representa a modificação ou estado do orgão \emph{externo} e o objecto \emph{externo}, que produziu essa modificação ou estado. A modificação ou o estado é percebido de um modo \emph{directo}, ao passo que o objecto é percebido de um modo \emph{indirecto}; porque, como dissemos, o objecto não influe no sentido \emph{interno} senão por meio da sensação \emph{externa}. --- D'ahi a differença entre a sensação \emph{externa} e a \emph{interna}. O objecto \emph{directo} da sensação \emph{externa} é o proprio ente \emph{externo}, que actuou \emph{directamente} nos sentidos \emph{externos}; mas o objecto \emph{directo} da sensação \emph{interna} é a \emph{paixão} dos sentidos \emph{externos}, isto é, é a propria sensação \emph{externa}, emquanto é modificação do sujeito sensiente.}.
